% MARC to TTL Mapping Documentation
% Hebrew Manuscripts Ontology Converter
% Compiled with XeLaTeX

\documentclass[12pt,a4paper]{article}

% Load packages BEFORE bidi (required order)
\usepackage{xcolor}
\usepackage[top=2.5cm, bottom=2.5cm, left=2.5cm, right=2.5cm]{geometry}
\usepackage{longtable}
\usepackage{array}
\usepackage{etoolbox}
\usepackage{hyperref}
\hypersetup{
    colorlinks=true,
    linkcolor=blue,
    urlcolor=blue,
    unicode=true
}

% Font and language setup
\usepackage{fontspec}
\usepackage{polyglossia}
\setmainlanguage{hebrew}
\setotherlanguage{english}

% Use FreeSerif font (available in Overleaf)
\setmainfont{FreeSerif}
\newfontfamily\hebrewfont{FreeSerif}
\newfontfamily\englishfont{FreeSerif}
\newfontfamily\hebrewfonttt{FreeMono}

% Load bidi LAST (after all other packages)
\usepackage{bidi}

% Custom column type for LTR English text (left-aligned)
% Use \leftskip=0pt to ensure left alignment in LTR mode
\newcolumntype{L}[1]{>{\setLTR\footnotesize\raggedright\arraybackslash}p{#1}<{\unsetLTR}}

% Small text column type for RTL Hebrew (right-aligned by default)
\newcolumntype{S}[1]{>{\footnotesize\raggedright\arraybackslash}p{#1}}

% Make all tables use smaller font by default
\AtBeginEnvironment{tabular}{\footnotesize}
\AtBeginEnvironment{longtable}{\footnotesize}

% Custom commands - use \texttt for monospace, LTR handled by column type L
\newcommand{\marcfield}[1]{\texttt{#1}}
\newcommand{\ttlprop}[1]{\texttt{#1}}
\newcommand{\ttlclass}[1]{\texttt{#1}}

% Hebrew parentheses - swap for correct RTL rendering
\newcommand{\heb}[1]{)#1(}

\definecolor{newaddition}{RGB}{0,128,0}
\newcommand{\new}[1]{\textcolor{newaddition}{#1}}
\newcommand{\newclass}[1]{\textcolor{newaddition}{\textbf{#1}}}
\newcommand{\newprop}[1]{\textcolor{newaddition}{#1}}

% Document
\begin{document}

% Title page
\begin{center}
{\LARGE\bfseries מיפוי MARC ל-TTL}\\[0.5cm]
{\Large מערכת ההמרה של אונטולוגיית כתבי יד עבריים}\\[0.5cm]
{\large גרסה 1.5 — דצמבר 2025}\\[0.3cm]
{\normalsize \new{כולל אינטגרציה עם מאגר הרשויות הלאומי — מזל}}\\[1cm]
\end{center}

\tableofcontents
\newpage

% =============================================================================
\section{מבוא}
% =============================================================================

מסמך זה מתאר את המיפוי בין שדות MARC לישויות ומאפיינים באונטולוגיית כתבי יד עבריים \textenglish{(TTL)}. מערכת ההמרה מעבדת רשומות MARC מהספרייה הלאומית של ישראל וממירה אותן לגרפים RDF בהתאם לאונטולוגיה בגרסה 1.5.

\new{החידוש המרכזי בגרסה זו הוא אינטגרציה עם מאגר הרשויות הלאומי מזל, המאפשר שימוש במזהי NLI רשמיים במקום URIs מקומיים.}

\subsection{מבנה המסמך}

\begin{itemize}
\item \textbf{פרק 2}: סקירת ארכיטקטורת המערכת
\item \textbf{פרק 3}: מיפוי שדות בקרה \textenglish{(00X)}
\item \textbf{פרק 4}: מיפוי שדות כותר ומחבר \textenglish{(1XX-2XX)}
\item \textbf{פרק 5}: מיפוי שדות פיזיים \textenglish{(3XX)}
\item \textbf{פרק 6}: מיפוי שדות הערות \textenglish{(5XX)}
\item \textbf{פרק 7}: מיפוי שדות נושא \textenglish{(6XX)}
\item \textbf{פרק 8}: מיפוי שדות נוספים \textenglish{(7XX-8XX)}
\item \textbf{פרק 9}: תכונות אונטולוגיה v1.4
\item \textbf{פרק 10}: אוצרות מילים מבוקרים
\end{itemize}

\subsection{עקרונות מנחים}

\begin{enumerate}
\item \textbf{שימור נתונים}: כל מידע מ-MARC נשמר ב-RDF
\item \textbf{סמנטיקה מפורשת}: כל נתון מקבל מבנה אונטולוגי מתאים
\item \textbf{URI עקביים}: יצירת מזהים עקביים לישויות
\item \textbf{מטא-נתונים אפיסטמולוגיים}: סימון רמת וודאות ומקור ייחוס
\end{enumerate}

\newpage
% =============================================================================
\section{ארכיטקטורת המערכת}
% =============================================================================

\subsection{רכיבי המערכת}

מערכת ההמרה מורכבת מהרכיבים הבאים:

\begin{longtable}{|L{5cm}|p{11cm}|}
\hline
\textbf{רכיב} & \textbf{תיאור} \\
\hline
\endfirsthead
\hline
\textbf{רכיב} & \textbf{תיאור} \\
\hline
\endhead
\hline
\endfoot
marc\_reader.py & קורא קבצי MARC בפורמט .mrc \\
\hline
field\_handlers.py & מטפלים ייעודיים לכל סוג שדה MARC \\
\hline
uri\_generator.py & יצירת URI עקביים לישויות \new{— כולל אינטגרציה עם מזל} \\
\hline
graph\_builder.py & בניית גרף RDF מהנתונים המופקים \\
\hline
vocabularies.py & מיפוי אוצרות מילים מבוקרים \\
\hline
namespaces.py & הגדרת מרחבי שמות RDF \\
\hline
\new{authority/mazal\_matcher.py} & \new{שירות התאמה לרשומות רשות מזל} \\
\hline
\new{authority/mazal\_index.py} & \new{אינדקס SQLite של 4.5M רשומות רשות NLI} \\
\hline
\end{longtable}

\subsection{תהליך ההמרה}

\begin{enumerate}
\item \textbf{קריאת MARC}: קבצי .mrc נקראים באמצעות pymarc
\item \textbf{פענוח שדות}: כל שדה מעובד על ידי handler מתאים
\item \textbf{חילוץ נתונים}: נתונים מאוחסנים ב-ExtractedData
\item \textbf{יצירת URI}: נוצרים מזהים ייחודיים לכל ישות
\item \textbf{בניית גרף}: נוצרות טריפלות RDF
\item \textbf{סריאליזציה}: הגרף נשמר בפורמט TTL
\end{enumerate}

\subsection{מרחבי שמות}

\begin{longtable}{|L{3cm}|L{13cm}|}
\hline
\textbf{קידומת} & \textbf{URI} \\
\hline
\endfirsthead
\hline
\textbf{קידומת} & \textbf{URI} \\
\hline
\endhead
\hline
\endfoot
hm: & http://www.ontology.org.il/HebrewManuscripts/2025-12-06\# \\
\hline
lrmoo: & http://iflastandards.info/ns/lrm/lrmoo/ \\
\hline
cidoc-crm: & http://www.cidoc-crm.org/cidoc-crm/ \\
\hline
rdf: & http://www.w3.org/1999/02/22-rdf-syntax-ns\# \\
\hline
rdfs: & http://www.w3.org/2000/01/rdf-schema\# \\
\hline
xsd: & http://www.w3.org/2001/XMLSchema\# \\
\hline
skos: & http://www.w3.org/2004/02/skos/core\# \\
\hline
owl: & http://www.w3.org/2002/07/owl\# \\
\hline
\new{nli:} & \new{https://www.nli.org.il/en/authorities/} \\
\hline
\end{longtable}

\subsection{\new{אינטגרציה עם מזל — מאגר הרשויות הלאומי}}

\new{החל מגרסה 1.5, מערכת ההמרה משתלבת עם מאגר הרשויות הישראלי מזל:}

\begin{itemize}
\item \new{אישים \textenglish{(MARC 100/700)} — התאמה לרשומות Person}
\item \new{מקומות \textenglish{(MARC 260/264)} — התאמה לרשומות Place}
\item \new{יצירות \textenglish{(MARC 130/240)} — התאמה לרשומות Work}
\item \new{גופים \textenglish{(MARC 110/710)} — התאמה לרשומות Corporate}
\end{itemize}

\new{התהליך:}
\begin{enumerate}
\item \new{השם מנורמל ומחופש באינדקס מזל}
\item \new{אם נמצאה התאמה — משתמשים ב-URI של NLI}
\item \new{אם לא נמצאה — משתמשים ב-URI מקומי \textenglish{(hm:Person\_...)}}
\end{enumerate}

\new{דוגמה:}
\begin{LTR}
\begin{verbatim}
# Input: MARC 100 $a "משה בן מימון" $d "1138-1204"

# Output with Mazal match:
<https://www.nli.org.il/en/authorities/987007388484005171>
    a cidoc-crm:E21_Person ;
    rdfs:label "משה בן מימון"@he ;
    hm:external_identifier_nli "987007388484005171" .
\end{verbatim}
\end{LTR}

\subsection{תבניות URI}

להלן הרשימה המלאה של תבניות URI עבור כל סוגי הישויות באונטולוגיה:

\begin{longtable}{|p{2cm}|L{4cm}|L{10cm}|}
\hline
\textbf{ישות} & \textbf{תבנית} & \textbf{דוגמה מלאה} \\
\hline
\endfirsthead
\hline
\textbf{ישות} & \textbf{תבנית} & \textbf{דוגמה מלאה} \\
\hline
\endhead
\hline
\endfoot
כתב יד & hm:MS\_\{id\} & hm:MS\_990012345680205171 \\
\hline
יצירה & hm:Work\_\{title\}[\_by\_\{author\}] & hm:Work\_משנה\_תורה\_by\_רמבם \\
\hline
ביטוי & hm:Expression\_\{title\}\_in\_\{id\} & hm:Expression\_משנה\_תורה\_in\_990012345680205171 \\
\hline
אדם & hm:Person\_\{name\} & hm:Person\_רמבם \\
\hline
ארגון & hm:Group\_\{name\} & hm:Group\_ישיבת\_וולוז׳ין \\
\hline
מקום & hm:Place\_\{name\} & hm:Place\_ספרד \\
\hline
פרק זמן & hm:TimeSpan\_\{label\} & hm:TimeSpan\_1407 \\
\hline
אירוע ייצור & hm:Production\_\{id\} & hm:Production\_990012345680205171 \\
\hline
אירוע יצירה & hm:WorkCreation\_\{title\} & hm:WorkCreation\_משנה\_תורה\_by\_רמבם \\
\hline
אירוע בעלות & hm:Ownership\_\{id\}\_\{seq\} & hm:Ownership\_990012345680205171\_1 \\
\hline
קטלוג & hm:Catalog\_\{name\} & hm:Catalog\_Institute\_of\_Hebrew\_Manuscripts \\
\hline
קולופון & hm:Colophon\_\{id\} & hm:Colophon\_990012345680205171 \\
\hline
נושא & hm:Subject\_\{term\} & hm:Subject\_הלכה \\
\hline
כריכה & hm:Binding\_\{id\} & hm:Binding\_990012345680205171 \\
\hline
יחידה ביב' & hm:BibUnit\_\{id\}\_\{seq\} & hm:BibUnit\_990012345680205171\_1 \\
\hline
יחידה קוד' & hm:CU\_\{id\}\_\{seq\} & hm:CU\_990012345680205171\_1 \\
\hline
יחידה פליאו' & hm:PU\_\{id\}\_\{seq\} & hm:PU\_990012345680205171\_1 \\
\hline
היררכיה קוד' & hm:Hierarchy\_\{id\} & hm:Hierarchy\_990012345680205171 \\
\hline
תצוגה קטלוגית & hm:CatalogingView\_\{id\} & hm:CatalogingView\_990012345680205171 \\
\hline
תצוגה פילולוגית & hm:PhilologicalView\_\{id\} & hm:PhilologicalView\_990012345680205171 \\
\hline
מסורת טקסט & hm:TextTradition\_\{name\} & hm:TextTradition\_משנה\_תורה \\
\hline
עד מסירה & hm:Witness\_\{title\}\_in\_\{id\} & hm:Witness\_משנה\_תורה\_in\_990012345680205171 \\
\hline
גשר פרדיגמות & hm:Bridge\_\{work\}\_to\_\{trad\} & hm:Bridge\_משנה\_תורה\_to\_TextTradition \\
\hline
גרסה טקסטואלית & hm:Variant\_\{id\}\_\{loc\} & hm:Variant\_990012345680205171\_15r\_L10 \\
\hline
התערבות סופר & hm:Intervention\_\{id\}\_\{seq\} & hm:Intervention\_990012345680205171\_1 \\
\hline
הפניה קאנונית & hm:CanonRef\_\{type\}\_\{ref\} & hm:CanonRef\_Bible\_Genesis\_1\_1 \\
\hline
מיקום טקסט & hm:Loc\_\{id\}\_\{folio\}[\_L\{n\}] & hm:Loc\_990012345680205171\_15r\_L10 \\
\hline
שרשרת ראיות & hm:EvidenceChain\_\{id\}\_\{field\} & hm:EvidenceChain\_990012345680205171\_date \\
\hline
סט רב-כרכי & hm:MultiVolumeSet\_\{name\} & hm:MultiVolumeSet\_משנה\_תורה \\
\hline
מבנה אנתולוגיה & hm:AnthologyStructure\_\{id\} & hm:AnthologyStructure\_990012345680205171 \\
\hline
\end{longtable}

\newpage
% =============================================================================
\section{מיפוי שדות בקרה \textenglish{(00X)}}
% =============================================================================

\subsection{שדה 001 — מספר בקרה}

\begin{longtable}{|L{4cm}|L{12cm}|}
\hline
\textbf{MARC} & \marcfield{001} \\
\hline
\textbf{תיאור} & מספר בקרה ייחודי של הרשומה \\
\hline
\textbf{מיפוי TTL} & \\
\hline
\textenglish{URI Pattern} & \ttlprop{hm:MS\_\{control\_number\}} \\
\hline
\textenglish{Property} & \ttlprop{hm:external\_identifier\_nli} \\
\hline
\textenglish{Datatype} & \ttlprop{xsd:string} \\
\hline
\end{longtable}

\textbf{דוגמה:}
\begin{LTR}
\begin{verbatim}
MARC: 001 990012345680205171
TTL:  hm:MS_990012345680205171 
        hm:external_identifier_nli "990012345680205171"^^xsd:string .
\end{verbatim}
\end{LTR}

\subsection{שדה 008 — נתונים קבועים}

שדה 008 מכיל מידע קבוע בפורמט ללא תת-שדות:

\begin{longtable}{|L{2cm}|L{2cm}|p{4cm}|L{8cm}|}
\hline
\textbf{מיקום} & \textbf{תווים} & \textbf{תיאור} & \textbf{מיפוי TTL} \\
\hline
\endfirsthead
\hline
\textbf{מיקום} & \textbf{תווים} & \textbf{תיאור} & \textbf{מיפוי TTL} \\
\hline
\endhead
\hline
\endfoot
\textenglish{06} & 1 & סוג תאריך MARC & קביעת אסטרטגיית תיארוך )ראו קודים למטה( \\
\hline
\textenglish{07-10} & 4 & תאריך התחלה & \ttlprop{hm:earliest\_possible\_date} )דרך \ttlprop{data.dates}( \\
\hline
\textenglish{11-14} & 4 & תאריך סיום & \ttlprop{hm:latest\_possible\_date} )דרך \ttlprop{data.dates}( \\
\hline
\textenglish{15-17} & 3 & קוד מקום & נחלץ אך לא נכתב לגרף )לשימוש עתידי( \\
\hline
\textenglish{35-37} & 3 & קוד שפה & \ttlprop{cidoc-crm:P72\_has\_language} \\
\hline
\end{longtable}

\textbf{קודי סוג תאריך MARC )מיקום 06(:}
\begin{itemize}
\item \textenglish{s} — תאריך בודד
\item \textenglish{q} — תאריך מסופק
\item \textenglish{m} — תאריכים מרובים
\item \textenglish{n} — לא ידוע
\end{itemize}

\textbf{פורמט תאריך )v1.4( — \ttlprop{hm:has\_date\_format}:}

המערכת מזהה ומסווגת תאריכים לפי 4 פורמטים:

\begin{longtable}{|p{3cm}|L{4.5cm}|L{4.5cm}|p{4cm}|}
\hline
\textbf{פורמט} & \textbf{ערך TTL} & \textbf{דוגמה} & \textbf{עברית} \\
\hline
\endfirsthead
\hline
\textbf{פורמט} & \textbf{ערך TTL} & \textbf{דוגמה} & \textbf{עברית} \\
\hline
\endhead
\hline
\endfoot
שנה גרגוריאנית & \ttlclass{hm:GregorianYear} & 1407, 1523 & שנה גרגוריאנית \\
\hline
שנה עברית & \ttlclass{hm:HebrewYear} & תשפ״ד, ה'קס״ז & שנה עברית \\
\hline
תאריך מלא & \ttlclass{hm:FullDate} & 15/03/1407 & תאריך מלא \\
\hline
טקסט חופשי & \ttlclass{hm:UnstructuredDate} & "המאה ה-15" & טקסט חופשי \\
\hline
\end{longtable}

\newpage
% =============================================================================
\section{מיפוי שדות כותר ומחבר \textenglish{(1XX-2XX)}}
% =============================================================================

\subsection{שדה 100 — מחבר ראשי}

\begin{longtable}{|L{2.5cm}|p{4.5cm}|L{9cm}|}
\hline
\textbf{תת-שדה} & \textbf{תיאור} & \textbf{מיפוי TTL} \\
\hline
\endfirsthead
\hline
\textbf{תת-שדה} & \textbf{תיאור} & \textbf{מיפוי TTL} \\
\hline
\endhead
\hline
\endfoot
\marcfield{\$a} & שם המחבר & \ttlprop{rdfs:label} על \ttlclass{cidoc-crm:E21\_Person} \\
\hline
\marcfield{\$d} & תאריכי חיים & \ttlprop{cidoc-crm:P82a\_begin\_of\_the\_begin} + \ttlprop{cidoc-crm:P82b\_end\_of\_the\_end} \\
\hline
\marcfield{\$e} & תפקיד & \ttlprop{cidoc-crm:P14\_carried\_out\_by} עם role \\
\hline
\marcfield{\$0} & מזהה סמכות & \ttlprop{hm:external\_uri\_nli} \\
\hline
\end{longtable}

\textbf{קשר לישויות:}
\begin{itemize}
\item יוצר \ttlclass{cidoc-crm:E21\_Person}
\item יוצר \ttlclass{lrmoo:F27\_Work\_Creation}
\item מקשר \ttlprop{lrmoo:R16\_created} ליצירה
\item מקשר \ttlprop{cidoc-crm:P14\_carried\_out\_by} לאדם
\end{itemize}

\textbf{דוגמה:}
\begin{LTR}
\begin{verbatim}
MARC: 100 1# $a רמב"ם $d 1138-1204
TTL:
hm:Person_רמבם a cidoc-crm:E21_Person ;
  rdfs:label "רמב\"ם"@he ;
  cidoc-crm:P82a_begin_of_the_begin 1138 ;
  cidoc-crm:P82b_end_of_the_end 1204 .
  
  hm:WorkCreation_... a lrmoo:F27_Work_Creation ;
  cidoc-crm:P14_carried_out_by hm:Person_רמבם ;
  lrmoo:R16_created hm:Work_... .
\end{verbatim}
\end{LTR}

\subsection{שדה 110 — ארגון ראשי}

\begin{longtable}{|L{2.5cm}|p{4.5cm}|L{9cm}|}
\hline
\textbf{תת-שדה} & \textbf{תיאור} & \textbf{מיפוי TTL} \\
\hline
\endfirsthead
\hline
\textbf{תת-שדה} & \textbf{תיאור} & \textbf{מיפוי TTL} \\
\hline
\endhead
\hline
\endfoot
\marcfield{\$a} & שם הארגון & \ttlprop{rdfs:label} על \ttlclass{cidoc-crm:E74\_Group} \\
\hline
\marcfield{\$0} & מזהה סמכות & \ttlprop{hm:external\_uri\_nli} \\
\hline
\end{longtable}

\subsection{שדה 245 — כותר}

\begin{longtable}{|L{2.5cm}|p{4.5cm}|L{9cm}|}
\hline
\textbf{תת-שדה} & \textbf{תיאור} & \textbf{מיפוי TTL} \\
\hline
\endfirsthead
\hline
\textbf{תת-שדה} & \textbf{תיאור} & \textbf{מיפוי TTL} \\
\hline
\endhead
\hline
\endfoot
\marcfield{\$a} & כותר ראשי & \ttlprop{hm:has\_title} + \ttlprop{rdfs:label} \\
\hline
\marcfield{\$b} & כותר משנה & מצורף לכותר המלא עם ":" \textenglish{(hm:has\_title)} \\
\hline
\marcfield{\$n} & מספר חלק & נחלץ אך לא נכתב לגרף )לשימוש עתידי( \\
\hline
\marcfield{\$p} & שם חלק & נחלץ אך לא נכתב לגרף )לשימוש עתידי( \\
\hline
\end{longtable}

\textbf{ישויות נוצרות:}
\begin{itemize}
\item \ttlclass{lrmoo:F1\_Work} — היצירה המופשטת
\item \ttlclass{lrmoo:F2\_Expression} — הביטוי הספציפי בכתב היד
\item קשר \ttlprop{lrmoo:R3\_is\_realised\_in} בין Expression ל-Work
\item קשר \ttlprop{lrmoo:R4\_embodies} בין Manuscript ל-Expression
\end{itemize}

\subsection{שדה 246 — כותרות חלופיות}

\begin{longtable}{|L{2.5cm}|p{4.5cm}|L{9cm}|}
\hline
\textbf{תת-שדה} & \textbf{תיאור} & \textbf{מיפוי TTL} \\
\hline
\marcfield{\$a} & כותר חלופי & \ttlprop{hm:has\_title} נוסף \\
\hline
\end{longtable}

\subsection{שדות 260/264 — מקום ותאריך}

\begin{longtable}{|L{2.5cm}|p{4.5cm}|L{9cm}|}
\hline
\textbf{תת-שדה} & \textbf{תיאור} & \textbf{מיפוי TTL} \\
\hline
\endfirsthead
\hline
\textbf{תת-שדה} & \textbf{תיאור} & \textbf{מיפוי TTL} \\
\hline
\endhead
\hline
\endfoot
\marcfield{\$a} & מקום & \ttlclass{cidoc-crm:E53\_Place} + \ttlprop{cidoc-crm:P7\_took\_place\_at} \\
\hline
\marcfield{\$c} & תאריך & \ttlclass{cidoc-crm:E52\_Time-Span} + \ttlprop{cidoc-crm:P4\_has\_time-span} \\
\hline
\end{longtable}

\textbf{עיבוד תאריכים \textenglish{(v1.4)} — סיווג פורמט:}

המערכת מזהה אוטומטית את פורמט התאריך:

\begin{enumerate}
\item \textbf{שנה גרגוריאנית} \textenglish{(hm:GregorianYear)}: זיהוי \textenglish{YYYY} )למשל: 1407, 1523(
\item \textbf{שנה עברית} \textenglish{(hm:HebrewYear)}: זיהוי תבניות עבריות )תשפ״ד, ה'קס״ז, שנת תש״ה(
\item \textbf{תאריך מלא} \textenglish{(hm:FullDate)}: זיהוי \textenglish{DD/MM/YYYY} או תאריך עברי מלא )כ״ג אדר תשפ״ד(
\item \textbf{טקסט חופשי} \textenglish{(hm:UnstructuredDate)}: כל מה שלא מתאים לפורמטים הנ״ל )למשל: "המאה ה-15"(
\end{enumerate}

\textbf{מאפיינים נוספים:}
\begin{itemize}
\item \ttlprop{hm:has\_date\_format} — מצביע לסוג הפורמט \textenglish{(hm:DateFormatType)}
\item \ttlprop{hm:date\_original\_string} — המחרוזת המקורית כפי שהופיעה במקור \textenglish{(xsd:string)}
\item סימוני אי-ודאות: \textenglish{[, ?, ca.}, בערך, לערך
\item המרה לטווח: \ttlprop{hm:earliest\_possible\_date}, \ttlprop{hm:latest\_possible\_date} \textenglish{(xsd:integer)}
\end{itemize}

\newpage
% =============================================================================
\section{מיפוי שדות פיזיים \textenglish{(3XX)}}
% =============================================================================

\subsection{שדה 300 — תיאור פיזי}

\begin{longtable}{|L{2.5cm}|p{4.5cm}|L{9cm}|}
\hline
\textbf{תת-שדה} & \textbf{תיאור} & \textbf{מיפוי TTL} \\
\hline
\endfirsthead
\hline
\textbf{תת-שדה} & \textbf{תיאור} & \textbf{מיפוי TTL} \\
\hline
\endhead
\hline
\endfoot
\marcfield{\$a} & היקף & \ttlprop{hm:has\_number\_of\_folios} \textenglish{(xsd:integer)} \\
\hline
\marcfield{\$b} & פרטים נוספים & נחלץ אך לא נכתב לגרף )לשימוש עתידי( \\
\hline
\marcfield{\$c} & מידות & \ttlprop{hm:has\_height\_mm}, \ttlprop{hm:has\_width\_mm} \\
\hline
\end{longtable}

\textbf{פענוח היקף:}

המערכת מזהה תבניות שונות:
\begin{itemize}
\item "248 leaves" → 248
\item "150 ff." → 150
\item "100 דפים" → 100
\item "[5], 120 leaves" → 125
\end{itemize}

\textbf{פענוח מידות:}

\begin{itemize}
\item "280 x 200 mm" → height=280, width=200
\item "28 cm" → height=280
\item "28 x 20 cm" → height=280, width=200
\end{itemize}

\subsection{שדה 340 — מדיום פיזי}

\begin{longtable}{|L{2.5cm}|p{4.5cm}|L{9cm}|}
\hline
\textbf{תת-שדה} & \textbf{תיאור} & \textbf{מיפוי TTL} \\
\hline
\marcfield{\$a} & חומר & \ttlprop{hm:has\_material} → \ttlclass{cidoc-crm:E57\_Material} \\
\hline
\end{longtable}

\textbf{מיפוי חומרים:}

\begin{longtable}{|p{4cm}|L{5cm}|L{7cm}|}
\hline
\textbf{מונח MARC} & \textbf{ערך TTL} & \textbf{עברית} \\
\hline
\endfirsthead
\hline
\textbf{מונח MARC} & \textbf{ערך TTL} & \textbf{עברית} \\
\hline
\endhead
\hline
\endfoot
parchment / vellum & Parchment & קלף \\
\hline
paper & Paper & נייר \\
\hline
papyrus & Papyrus & פפירוס \\
\hline
\end{longtable}

\newpage
% =============================================================================
\section{מיפוי שדות הערות \textenglish{(5XX)}}
% =============================================================================

\subsection{שדה 500 — הערות כלליות}

שדה 500 הוא המורכב ביותר לעיבוד כיוון שהוא מכיל מידע חופשי. המערכת מבצעת ניתוח טקסט לזיהוי:

\subsubsection{זיהוי סוג כתב}

\begin{longtable}{|p{4cm}|L{5cm}|L{7cm}|}
\hline
\textbf{מונח} & \textbf{ערך TTL} & \textbf{מאפיין} \\
\hline
\endfirsthead
\hline
\textbf{מונח} & \textbf{ערך TTL} & \textbf{מאפיין} \\
\hline
\endhead
\hline
\endfoot
אשכנזי / Ashkenazi & AshkenaziScript & \ttlprop{hm:has\_script\_type} \\
\hline
ספרדי / Sephardic & SepharadicScript & \ttlprop{hm:has\_script\_type} \\
\hline
איטלקי / Italian & ItalianScript & \ttlprop{hm:has\_script\_type} \\
\hline
ביזנטי / Byzantine & ByzantineScript & \ttlprop{hm:has\_script\_type} \\
\hline
תימני / Yemenite & YemeniteScript & \ttlprop{hm:has\_script\_type} \\
\hline
מזרחי / Oriental & OrientalScript & \ttlprop{hm:has\_script\_type} \\
\hline
פרסי / Persian & PersianScript & \ttlprop{hm:has\_script\_type} \\
\hline
\end{longtable}

\subsubsection{זיהוי מצב כתב}

\begin{longtable}{|p{4cm}|L{5cm}|L{7cm}|}
\hline
\textbf{מונח} & \textbf{ערך TTL} & \textbf{מאפיין} \\
\hline
\endfirsthead
\hline
\textbf{מונח} & \textbf{ערך TTL} & \textbf{מאפיין} \\
\hline
\endhead
\hline
\endfoot
מרובע / Square & SquareScript & \ttlprop{hm:has\_Script\_Mode} \\
\hline
בינוני / Semi-cursive & SemicursiveScript & \ttlprop{hm:has\_Script\_Mode} \\
\hline
רהוט / Cursive & CursiveScript & \ttlprop{hm:has\_Script\_Mode} \\
\hline
\end{longtable}

\subsubsection{זיהוי קולופון}

מילות מפתח: קולופון, colophon, כתב הסופר, כתב המעתיק

\begin{LTR}
\begin{verbatim}
TTL:
  hm:Colophon_{control_number} a hm:Colophon ;
    hm:colophon_text "טקסט הקולופון"@he .
  
  hm:MS_{control_number} hm:has_colophon 
    hm:Colophon_{control_number} .
\end{verbatim}
\end{LTR}

\subsubsection{זיהוי יחידות קודיקולוגיות )v1.4(}

המערכת מזהה תבניות המעידות על מבנה מורכב ושומרת אותם ב-\textenglish{ExtractedData.codicological\_units}.

\textbf{הערה חשובה:} נתונים אלה נחלצים לצורך שימוש עתידי, אך לא נכתבים אוטומטית לגרף RDF. הוספת יחידות קודיקולוגיות לגרף מתבצעת דרך API הייעודי \textenglish{(add\_codicological\_unit)}.

תבניות מזוהות:
\begin{itemize}
\item "יחידה 1", "יחידה 2" — מספור יחידות
\item "חלק 1", "part 1" — חלקים
\item "קוויון 1-4" — טווח קוויונים
\item "ff. 1r-50v" — טווח דפים
\item אנתולוגיה, anthology, sammelband — אינדיקטורים של קובץ
\end{itemize}

\subsubsection{זיהוי התערבויות סופר )v1.4(}

המערכת מזהה אינדיקטורים להתערבויות סופר מטקסט חופשי ושומרת אותם ב-\textenglish{ExtractedData.scribal\_interventions}. 

\textbf{הערה חשובה:} נתונים אלה נחלצים לצורך שימוש עתידי, אך לא נכתבים אוטומטית לגרף RDF. הוספת התערבויות סופר לגרף מתבצעת דרך API הייעודי \textenglish{(add\_scribal\_intervention)}.

\begin{longtable}{|p{4cm}|L{5cm}|L{7cm}|}
\hline
\textbf{מונח} & \textbf{סוג התערבות} & \textbf{TTL Instance} \\
\hline
\endfirsthead
\hline
\textbf{מונח} & \textbf{סוג התערבות} & \textbf{TTL Instance} \\
\hline
\endhead
\hline
\endfoot
correction / תיקון & תיקון & \ttlclass{hm:Correction\_type} \\
\hline
erasure / מחיקה & מחיקה & \ttlclass{hm:Erasure\_type} \\
\hline
interlinear\_addition / הוספה בין השורות & הוספה בין השורות & \ttlclass{hm:Interlinear\_addition\_type} \\
\hline
marginal\_gloss / ביאור בשוליים & ביאור בשוליים & \ttlclass{hm:Marginal\_gloss\_type} \\
\hline
later\_hand / יד מאוחרת & יד מאוחרת & \ttlclass{hm:Later\_hand\_type} \\
\hline
\end{longtable}

\textbf{חשיבות גרסאות טקסטואליות:}

\begin{longtable}{|p{4cm}|L{5cm}|L{7cm}|}
\hline
\textbf{חשיבות} & \textbf{תיאור} & \textbf{TTL Instance} \\
\hline
\endfirsthead
\hline
\textbf{חשיבות} & \textbf{תיאור} & \textbf{TTL Instance} \\
\hline
\endhead
\hline
\endfoot
orthographic / גרסה כתיבית & הבדלי כתיב בלבד & \ttlclass{hm:Orthographic\_variant} \\
\hline
lexical / גרסה לקסיקלית & בחירת מילה שונה & \ttlclass{hm:Lexical\_variant} \\
\hline
syntactic / גרסה תחבירית & סדר מילים שונה & \ttlclass{hm:Syntactic\_variant} \\
\hline
semantic / גרסה משמעותית & שינוי משמעות & \ttlclass{hm:Semantic\_variant} \\
\hline
addition / הוספה & טקסט קיים בעד זה אך לא בסטנדרטי & \ttlclass{hm:Addition\_variant} \\
\hline
omission / השמטה & טקסט חסר בעד זה & \ttlclass{hm:Omission\_variant} \\
\hline
\end{longtable}

\subsubsection{זיהוי הפניות קאנוניות )v1.4(}

המערכת מזהה הפניות לטקסטים קאנוניים ושומרת אותם ב-\textenglish{ExtractedData.canonical\_references}.

\textbf{הערה חשובה:} נתונים אלה נחלצים לצורך שימוש עתידי, אך לא נכתבים אוטומטית לגרף RDF. הוספת הפניות קאנוניות לגרף מתבצעת דרך API הייעודי \textenglish{(add\_canonical\_reference)}.

\textbf{סטטוס יחידה קודיקולוגית )תפקיד היחידה בכתב היד(:}

\begin{longtable}{|p{4cm}|p{6cm}|L{6cm}|}
\hline
\textbf{סטטוס} & \textbf{תיאור} & \textbf{TTL Instance} \\
\hline
\endfirsthead
\hline
\textbf{סטטוס} & \textbf{תיאור} & \textbf{TTL Instance} \\
\hline
\endhead
\hline
\endfoot
\textenglish{CoreUnit} & יחידת ליבה - חלק מהתפיסה המקורית & \ttlclass{hm:CoreUnit\_status} \\
\hline
\textenglish{LaterAddition} & תוספת מאוחרת - נוספה לאחר היצירה & \ttlclass{hm:LaterAddition\_status} \\
\hline
\textenglish{BinderAddition} & תוספת כורך - נוספה בתהליך הכריכה & \ttlclass{hm:BinderAddition\_status} \\
\hline
\textenglish{UnrelatedFragment} & שבר לא קשור לתוכן & \ttlclass{hm:UnrelatedFragment\_status} \\
\hline
\textenglish{ProtectiveLeaf} & דף מגן )דף פורזץ( & \ttlclass{hm:ProtectiveLeaf\_status} \\
\hline
\end{longtable}

\textbf{מצב פיזי \textenglish{(ConditionType)}:}

\begin{longtable}{|p{4cm}|p{6cm}|L{6cm}|}
\hline
\textbf{מצב} & \textbf{תיאור} & \textbf{TTL Instance} \\
\hline
\endfirsthead
\hline
\textbf{מצב} & \textbf{תיאור} & \textbf{TTL Instance} \\
\hline
\endhead
\hline
\endfoot
\textenglish{Excellent} & מצב מצוין & \ttlclass{hm:Excellent} \\
\hline
\textenglish{Good} & מצב טוב & \ttlclass{hm:Good} \\
\hline
\textenglish{Fair} & מצב סביר & \ttlclass{hm:Fair} \\
\hline
\textenglish{Poor} & מצב גרוע & \ttlclass{hm:Poor} \\
\hline
\textenglish{Fragmentary} & חלקי/קטעי & \ttlclass{hm:Fragmentary} \\
\hline
\end{longtable}

\textbf{ספרי תנ"ך:}
בראשית, שמות, ויקרא, במדבר, דברים, יהושע, שופטים, שמואל, מלכים, ישעיה, ירמיה, יחזקאל, תהלים, משלי, איוב...

\textbf{מסכתות תלמוד:}
ברכות, שבת, פסחים, יומא, סוכה, ביצה, ראש השנה, תענית, מגילה, מועד קטן, חגיגה, יבמות, כתובות, נדרים, נזיר, סוטה, גיטין, קידושין, בבא קמא, בבא מציעא, בבא בתרא, סנהדרין...

\subsubsection{שדות ExtractedData עבור v1.4}

המחלקה \textenglish{ExtractedData} מכילה שדות ייעודיים לתכונות v1.4:

\begin{longtable}{|L{5cm}|p{11cm}|}
\hline
\textbf{שדה} & \textbf{תיאור} \\
\hline
\endfirsthead
\hline
\textbf{שדה} & \textbf{תיאור} \\
\hline
\endhead
\hline
\endfoot
\textenglish{certainty\_levels} & מיפוי שם שדה → רמת וודאות \\
\hline
\textenglish{attribution\_sources} & מיפוי שם שדה → מקור ייחוס \\
\hline
\textenglish{data\_from\_colophon} & רשימת שדות שנגזרו מקולופון \\
\hline
\textenglish{codicological\_units} & רשימת יחידות קודיקולוגיות שזוהו \\
\hline
\textenglish{hierarchy\_type} & סוג היררכיה \textenglish{(Simple, Nested, Complex, Fragmentary)} \\
\hline
\textenglish{is\_anthology} & האם כתב יד מסוג אנתולוגיה \\
\hline
\textenglish{is\_multi\_volume} & האם כתב יד מסוג רב-כרכי \\
\hline
\textenglish{volume\_number} & מספר כרך )עבור כתבי יד רב-כרכיים( \\
\hline
\textenglish{scribal\_interventions} & רשימת התערבויות סופר שזוהו \\
\hline
\textenglish{canonical\_references} & רשימת הפניות קאנוניות שזוהו \\
\hline
\textenglish{text\_traditions} & רשימת מסורות טקסטואליות \\
\hline
\textenglish{textual\_variants} & רשימת גרסאות טקסטואליות \\
\hline
\end{longtable}

\subsection{שדה 505 — תוכן עניינים}

\begin{longtable}{|L{2.5cm}|p{4.5cm}|L{9cm}|}
\hline
\textbf{תת-שדה} & \textbf{תיאור} & \textbf{מיפוי TTL} \\
\hline
\endfirsthead
\hline
\textbf{תת-שדה} & \textbf{תיאור} & \textbf{מיפוי TTL} \\
\hline
\endhead
\hline
\endfoot
\marcfield{\$a} & תוכן מעוצב & פירוק לפריטים נפרדים \\
\hline
\marcfield{\$t} & כותר & יצירת \ttlclass{lrmoo:F1\_Work} + \ttlclass{lrmoo:F2\_Expression} \\
\hline
\marcfield{\$r} & אחריות & נחלץ אך לא נכתב לגרף )לשימוש עתידי( \\
\hline
\marcfield{\$g} & מידע נוסף & \ttlprop{hm:has\_folio\_range} \\
\hline
\end{longtable}

כל פריט תוכן יוצר:
\begin{itemize}
\item \ttlclass{lrmoo:F1\_Work} עבור היצירה
\item \ttlclass{lrmoo:F2\_Expression} עבור הביטוי בכתב היד
\item קשר \ttlprop{lrmoo:R4\_embodies} מכתב היד לביטוי
\end{itemize}

\subsection{שדה 510 — הפניות קטלוגיות}

\begin{longtable}{|L{2.5cm}|p{4.5cm}|L{9cm}|}
\hline
\textbf{תת-שדה} & \textbf{תיאור} & \textbf{מיפוי TTL} \\
\hline
\marcfield{\$a} & שם הקטלוג & \ttlclass{lrmoo:F3\_Manifestation} + \ttlprop{rdfs:label} \\
\hline
\marcfield{\$c} & מיקום & נחלץ אך לא נכתב לגרף )לשימוש עתידי( \\
\hline
\end{longtable}

קשרים:
\begin{itemize}
\item \ttlprop{cidoc-crm:P70i\_is\_documented\_in} — כתב היד מתועד בקטלוג
\item \ttlprop{cidoc-crm:P70\_documents} — הקטלוג מתעד את כתב היד
\end{itemize}

\subsection{שדה 561 — פרובננס}

\begin{longtable}{|L{2.5cm}|p{4.5cm}|L{9cm}|}
\hline
\textbf{תת-שדה} & \textbf{תיאור} & \textbf{מיפוי TTL} \\
\hline
\marcfield{\$a} & היסטוריית בעלות & \ttlprop{hm:ownership\_history} \\
\hline
\end{longtable}

\subsection{שדה 563 — כריכה}

\begin{longtable}{|L{2.5cm}|p{4.5cm}|L{9cm}|}
\hline
\textbf{תת-שדה} & \textbf{תיאור} & \textbf{מיפוי TTL} \\
\hline
\marcfield{\$a} & תיאור כריכה & \ttlclass{hm:Binding} + \ttlprop{rdfs:comment} \\
\hline
\end{longtable}

\newpage
% =============================================================================
\section{מיפוי שדות נושא \textenglish{(6XX)}}
% =============================================================================

\subsection{שדות 600-655 — נושאים}

\begin{longtable}{|L{2cm}|p{4cm}|L{4cm}|L{6cm}|}
\hline
\textbf{שדה} & \textbf{סוג} & \textbf{TTL Class} & \textbf{קשר} \\
\hline
\endfirsthead
\hline
\textbf{שדה} & \textbf{סוג} & \textbf{TTL Class} & \textbf{קשר} \\
\hline
\endhead
\hline
\endfoot
600 & אדם & \ttlclass{cidoc-crm:E21\_Person} & \ttlprop{cidoc-crm:P129\_is\_about} \\
\hline
610 & ארגון & \ttlclass{cidoc-crm:E74\_Group} & \ttlprop{cidoc-crm:P129\_is\_about} \\
\hline
650 & נושא & \ttlclass{hm:SubjectType} & \ttlprop{cidoc-crm:P129\_is\_about} \\
\hline
651 & מקום & \ttlclass{cidoc-crm:E53\_Place} & \ttlprop{cidoc-crm:P129\_is\_about} \\
\hline
655 & ז'אנר & \ttlclass{hm:SubjectType} & \ttlprop{cidoc-crm:P2\_has\_type} \\
\hline
\end{longtable}

\textbf{תת-שדות משותפים:}
\begin{itemize}
\item \marcfield{\$a} — מונח \textenglish{(rdfs:label)}
\item \marcfield{\$d} — תאריכים )לאנשים(
\item \marcfield{\$0} — מזהה סמכות \textenglish{(hm:external\_uri\_nli)}
\item \marcfield{\$2} — מקור אוצר מילים )נחלץ אך לא נכתב לגרף(
\end{itemize}

\newpage
% =============================================================================
\section{מיפוי שדות נוספים \textenglish{(7XX-8XX)}}
% =============================================================================

\subsection{שדות 700/710 — שמות נוספים}

\begin{longtable}{|L{2.5cm}|p{4.5cm}|L{9cm}|}
\hline
\textbf{תת-שדה} & \textbf{תיאור} & \textbf{מיפוי TTL} \\
\hline
\endfirsthead
\hline
\textbf{תת-שדה} & \textbf{תיאור} & \textbf{מיפוי TTL} \\
\hline
\endhead
\hline
\endfoot
\marcfield{\$a} & שם & \ttlclass{cidoc-crm:E21\_Person} / \ttlclass{cidoc-crm:E74\_Group} \\
\hline
\marcfield{\$d} & תאריכים & כמו שדה 100 \\
\hline
\marcfield{\$e} & תפקיד & קביעת סוג הקשר )כמו שדה 100( \\
\hline
\marcfield{\$4} & קוד תפקיד & נחלץ אך לא נכתב לגרף )לשימוש עתידי( \\
\hline
\end{longtable}

\textbf{מיפוי תפקידים:}

\begin{longtable}{|p{4cm}|L{5cm}|L{7cm}|}
\hline
\textbf{מונח} & \textbf{ערך} & \textbf{קשר TTL} \\
\hline
\endfirsthead
\hline
\textbf{מונח} & \textbf{ערך} & \textbf{קשר TTL} \\
\hline
\endhead
\hline
\endfoot
סופר / scribe / copyist & scribe & \ttlprop{cidoc-crm:P14\_carried\_out\_by} \textenglish{(Production)} \\
\hline
מאייר / illuminator & illuminator & contributor \\
\hline
פרשן / commentator & commentator & contributor \\
\hline
מתרגם / translator & translator & contributor \\
\hline
בעלים קודם / former owner & former\_owner & provenance \\
\hline
עורך / editor & editor & contributor \\
\hline
מלקט / compiler & compiler & contributor \\
\hline
\end{longtable}

\subsection{שדה 856 — מיקום אלקטרוני}

\begin{longtable}{|L{2.5cm}|p{4.5cm}|L{9cm}|}
\hline
\textbf{תת-שדה} & \textbf{תיאור} & \textbf{מיפוי TTL} \\
\hline
\marcfield{\$u} & URL & \ttlprop{hm:has\_digital\_representation\_url} \\
\hline
\marcfield{\$y} & טקסט קישור & נחלץ אך לא נכתב לגרף )לשימוש עתידי( \\
\hline
\marcfield{\$z} & הערה & נחלץ אך לא נכתב לגרף )לשימוש עתידי( \\
\hline
\end{longtable}

\newpage
% =============================================================================
\section{תכונות אונטולוגיה v1.4}
% =============================================================================

\subsection{מסגרת אפיסטמולוגית}

גרסה 1.4 מוסיפה מטא-נתונים לכל פריט מידע:

\subsubsection{רמות וודאות}

\begin{longtable}{|p{3cm}|L{5cm}|L{8cm}|}
\hline
\textbf{רמה} & \textbf{ערך TTL} & \textbf{שימוש} \\
\hline
\endfirsthead
\hline
\textbf{רמה} & \textbf{ערך TTL} & \textbf{שימוש} \\
\hline
\endhead
\hline
\endfoot
וודאי & \ttlprop{hm:Certain} & מדידות פיזיות, תעתיק ישיר \\
\hline
סביר & \ttlprop{hm:Probable} & רוב ייחוסי קטלוג \\
\hline
אפשרי & \ttlprop{hm:Possible} & תאריכים משוערים \\
\hline
לא וודאי & \ttlprop{hm:Uncertain} & נתונים מסופקים \\
\hline
\end{longtable}

\textbf{זיהוי אוטומטי:}
\begin{itemize}
\item סימני [, ?, ca. → \ttlprop{hm:Possible}
\item בערך, לערך → \ttlprop{hm:Possible}
\item תיארוך לפי מאה → \ttlprop{hm:Possible}
\item מדידות פיזיות → \ttlprop{hm:Certain}
\end{itemize}

\subsubsection{מקורות ייחוס}

\begin{longtable}{|p{4cm}|L{5cm}|L{7cm}|}
\hline
\textbf{מקור} & \textbf{ערך TTL} & \textbf{תיאור} \\
\hline
\endfirsthead
\hline
\textbf{מקור} & \textbf{ערך TTL} & \textbf{תיאור} \\
\hline
\endhead
\hline
\endfoot
קטלוג & \ttlprop{hm:CatalogAttribution} & ברירת מחדל לנתוני MARC \\
\hline
קולופון & \ttlprop{hm:ColophonAttribution} & מידע מקולופון \\
\hline
פליאוגרפי & \ttlprop{hm:PaleographicAttribution} & זיהוי סוג כתב \\
\hline
מומחה & \ttlprop{hm:ExpertAttribution} & ייחוס מומחה \\
\hline
AI & \ttlprop{hm:AIAttribution} & ייחוס חישובי \\
\hline
\end{longtable}

\subsubsection{סטטוסים אפיסטמולוגיים}

\begin{longtable}{|p{4cm}|L{5cm}|L{7cm}|}
\hline
\textbf{סטטוס} & \textbf{ערך TTL} & \textbf{תיאור} \\
\hline
\endfirsthead
\hline
\textbf{סטטוס} & \textbf{ערך TTL} & \textbf{תיאור} \\
\hline
\endhead
\hline
\endfoot
תצפית ישירה & \ttlprop{hm:DirectObservation} & מדידות פיסיות, תעתיק \\
\hline
פרשנות חוקר & \ttlprop{hm:ScholarlyInterpretation} & ייחוסים מפורשים \\
\hline
גזירה חישובית & \ttlprop{hm:ComputationalDerivation} & ייחוס AI \\
\hline
ירושה מקטלוג & \ttlprop{hm:CatalogInherited} & נתוני MARC לא מאומתים \\
\hline
ייחוס מסורתי & \ttlprop{hm:TraditionalAttribution} & ייחוס ללא אימות מודרני \\
\hline
\end{longtable}

\subsubsection{קטגוריות נתונים}

\begin{longtable}{|p{4cm}|L{5cm}|L{7cm}|}
\hline
\textbf{קטגוריה} & \textbf{ערך TTL} & \textbf{עובדתי?} \\
\hline
\endfirsthead
\hline
\textbf{קטגוריה} & \textbf{ערך TTL} & \textbf{עובדתי?} \\
\hline
\endhead
\hline
\endfoot
מדידה פיזית & PhysicalMeasurement & כן \\
\hline
תוכן קולופון & ColophonContent & כן \\
\hline
תצפית חומר & MaterialObservation & כן \\
\hline
ייחוס תאריך & DateAttribution & לא \\
\hline
ייחוס מחברות & AuthorshipAttribution & לא \\
\hline
זיהוי סופר & ScribeIdentification & לא \\
\hline
ייחוס מקום & PlaceAttribution & לא \\
\hline
זיהוי יצירה & WorkIdentification & לא \\
\hline
יחס טקסטואלי & TextualRelationship & לא \\
\hline
\end{longtable}

\subsection{פרדיגמה כפולה}

המערכת תומכת בשתי גישות תיאור:

\subsubsection{גישה קטלוגית )ביבליוגרפית(}

\begin{LTR}
\begin{verbatim}
TTL:
  hm:CatalogingView_{control_number} a hm:CatalogingView ;
    rdfs:label "Cataloging view for MS ..." ;
    hm:cataloging_work hm:Work_... ;
    hm:cataloging_expression hm:Expression_... ;
    hm:is_primary_paradigm true ;
    hm:paradigm_note "Data derived from NLI catalog..." .
  
  hm:MS_{control_number} hm:has_cataloging_perspective
    hm:CatalogingView_{control_number} .
\end{verbatim}
\end{LTR}

\subsubsection{גישה פילולוגית}

\begin{LTR}
\begin{verbatim}
TTL:
  hm:PhilologicalView_{control_number} a hm:PhilologicalView ;
    hm:view_type hm:PhilologicalParadigm ;
    hm:is_primary_paradigm false ;
    hm:paradigm_note "New Philology view: manuscript as 
                      unique cultural event" .
  
  hm:MS_{control_number} hm:has_philological_perspective
    hm:PhilologicalView_{control_number} .
\end{verbatim}
\end{LTR}

\subsubsection{גשר פרדיגמות}

\begin{LTR}
\begin{verbatim}
TTL:
  hm:Bridge_{work}_to_{tradition} a hm:ParadigmBridge ;
    hm:has_linked_work hm:Work_... ;
    hm:has_linked_tradition hm:TextTradition_... ;
    hm:paradigm_justification "..." .
\end{verbatim}
\end{LTR}

\subsubsection{מסורת טקסטואלית ועד מסירה}

מסורת טקסטואלית \textenglish{(TextTradition)} מייצגת את השתלשלות הטקסט בהעתקות. עד המסירה \textenglish{(TransmissionWitness)} מייצג את הכתב יד כאירוע תרבותי ייחודי בשרשרת המסירה.

\begin{LTR}
\begin{verbatim}
TTL:
  # Text Tradition
  hm:TextTradition_{name} a hm:TextTradition ;
    rdfs:label "Tradition of {name}" ;
    hm:tradition_name "{name}" .
  
  # Transmission Witness
  hm:Witness_{title}_in_{control_number} a hm:TransmissionWitness ;
    rdfs:label "Witness of {title} in MS {control_number}" ;
    hm:witnesses hm:TextTradition_{name} .
  
  # Direct link from manuscript to witness (shortcut property)
  hm:MS_{control_number} hm:has_philological_witness
    hm:Witness_{title}_in_{control_number} .
  
  # Full chain via PhilologicalView
  hm:PhilologicalView_{control_number} hm:philological_witness
    hm:Witness_{title}_in_{control_number} .
\end{verbatim}
\end{LTR}

הקשרים העיקריים:
\begin{itemize}
\item \ttlprop{hm:has\_philological\_witness} — קישור ישיר מכתב יד לעד מסירה
\item \ttlprop{hm:philological\_witness} — קישור מ-PhilologicalView לעד מסירה
\item \ttlprop{hm:witnesses} — קישור מעד מסירה למסורת טקסטואלית
\item \ttlprop{hm:tradition\_includes} — קישור הפוך — ממסורת לעד
\end{itemize}

\subsection{מבנה קודיקולוגי מקונן}

גרסה 1.4 תומכת ביחידות קודיקולוגיות מקוננות.

\subsubsection{סוגי היררכיה}

\begin{longtable}{|p{4cm}|L{5cm}|L{7cm}|}
\hline
\textbf{סוג} & \textbf{TTL Instance} & \textbf{תיאור} \\
\hline
\endfirsthead
\hline
\textbf{סוג} & \textbf{TTL Instance} & \textbf{תיאור} \\
\hline
\endhead
\hline
\endfoot
\textenglish{Simple} & \ttlclass{hm:SimpleHierarchy} & מבנה פשוט ללא קינון \\
\hline
\textenglish{Nested} & \ttlclass{hm:NestedHierarchy} & יחידות מקוננות \\
\hline
\textenglish{Complex} & \ttlclass{hm:ComplexHierarchy} & מבנה מורכב )כולל אנתולוגיה( \\
\hline
\textenglish{Fragmentary} & \ttlclass{hm:FragmentaryHierarchy} & מבנה קטוע/חסר \\
\hline
\end{longtable}

\subsubsection{דוגמת TTL}

\begin{LTR}
\begin{verbatim}
TTL:
  hm:Hierarchy_{control_number} a hm:CodicologicalHierarchy ;
    hm:hierarchy_type hm:NestedHierarchy ;
    hm:max_nesting_depth 2 .
  
  hm:MS_{control_number} hm:has_hierarchy 
    hm:Hierarchy_{control_number} .
  
  hm:CU_{control_number}_1 a hm:CompositeCodicologicalUnit ;
    hm:nesting_level 0 ;
    hm:unit_sequence 1 ;
    hm:has_sub_unit hm:CU_{control_number}_1_1 .
  
  hm:CU_{control_number}_1_1 a hm:AtomicCodicologicalUnit ;
    hm:is_sub_unit_of hm:CU_{control_number}_1 ;
    hm:nesting_level 1 ;
    hm:has_folio_range "1r-50v" ;
    hm:has_unit_status hm:CoreUnit_status .
\end{verbatim}
\end{LTR}

\subsection{שרשרת ראיות}

\begin{LTR}
\begin{verbatim}
TTL:
  hm:EvidenceChain_{control_number} a hm:EvidenceChain ;
    rdfs:label "Evidence chain for MS ..." .
  
  hm:MS_{control_number} hm:has_evidence_chain
    hm:EvidenceChain_{control_number} .
  
  hm:CatalogStep_{control_number} a hm:EvidenceStep ;
    hm:reasoning_text "Data imported from NLI MARC catalog" .
  
  hm:EvidenceChain_{control_number} hm:evidence_step
    hm:CatalogStep_{control_number} .
  
  hm:MS_{control_number} 
    hm:attribution_source hm:CatalogAttribution ;
    hm:has_epistemological_status hm:CatalogInherited ;
    hm:is_factual false .
\end{verbatim}
\end{LTR}

\subsection{API לתכונות מתקדמות}

המערכת מספקת API לתכונות v1.4 שאינן נכתבות אוטומטית מנתוני MARC:

\begin{longtable}{|L{5cm}|p{11cm}|}
\hline
\textbf{פונקציה} & \textbf{תיאור} \\
\hline
\endfirsthead
\hline
\textbf{פונקציה} & \textbf{תיאור} \\
\hline
\endhead
\hline
\endfoot
\textenglish{add\_codicological\_hierarchy} & יצירת מבנה היררכי לכתב יד מורכב \\
\hline
\textenglish{add\_codicological\_unit} & הוספת יחידה קודיקולוגית )תומך קינון( \\
\hline
\textenglish{add\_philological\_view} & הוספת פרספקטיבה פילולוגית \\
\hline
\textenglish{add\_text\_tradition} & יצירת מסורת טקסטואלית \\
\hline
\textenglish{add\_transmission\_witness} & קישור כתב יד למסורת טקסטואלית \\
\hline
\textenglish{add\_paradigm\_bridge} & גשר בין יצירה \textenglish{(Work)} למסורת \\
\hline
\textenglish{add\_textual\_variant} & הוספת גרסה טקסטואלית \\
\hline
\textenglish{add\_scribal\_intervention} & הוספת התערבות סופר \\
\hline
\textenglish{add\_canonical\_reference} & הוספת הפניה קאנונית \\
\hline
\textenglish{add\_detailed\_evidence\_chain} & שרשרת ראיות מפורטת \\
\hline
\end{longtable}

\newpage
% =============================================================================
\section{אוצרות מילים מבוקרים}
% =============================================================================

\subsection{קודי שפה}

\begin{longtable}{|p{3cm}|L{5cm}|L{8cm}|}
\hline
\textbf{קוד MARC} & \textbf{ערך TTL} & \textbf{עברית} \\
\hline
\endfirsthead
\hline
\textbf{קוד MARC} & \textbf{ערך TTL} & \textbf{עברית} \\
\hline
\endhead
\hline
\endfoot
heb & Hebrew & עברית \\
\hline
ara & Arabic & ערבית \\
\hline
arc & Aramaic & ארמית \\
\hline
jrb & JudeoArabic & יהודית-ערבית \\
\hline
lad & Ladino & לאדינו \\
\hline
yid & Yiddish & יידיש \\
\hline
jpr & JudeoPersian & יהודית-פרסית \\
\hline
lat & Latin & לטינית \\
\hline
grc & Greek & יוונית \\
\hline
per & Persian & פרסית \\
\hline
ger & German & גרמנית \\
\hline
spa & Spanish & ספרדית \\
\hline
ita & Italian & איטלקית \\
\hline
fre & French & צרפתית \\
\hline
eng & English & אנגלית \\
\hline
por & Portuguese & פורטוגזית \\
\hline
tur & Turkish & טורקית \\
\hline
und & Undetermined & לא מוגדר \\
\hline
mul & Multiple & מספר שפות \\
\hline
\end{longtable}

\textbf{הערה:} קודי שפה נוספים נתמכים על פי תקן ISO 639-2. הממיר מזהה אוטומטית קודים נפוצים ומתרגם אותם לשם השפה באנגלית.

\subsection{סוגי כתב}

\begin{longtable}{|p{4cm}|L{5cm}|L{7cm}|}
\hline
\textbf{מונחים} & \textbf{ערך TTL} & \textbf{URI} \\
\hline
\endfirsthead
\hline
\textbf{מונחים} & \textbf{ערך TTL} & \textbf{URI} \\
\hline
\endhead
\hline
\endfoot
אשכנזי / Ashkenazi & AshkenaziScript & \ttlprop{hm:AshkenaziScript} \\
\hline
ספרדי / Sephardic & SepharadicScript & \ttlprop{hm:SepharadicScript} \\
\hline
איטלקי / Italian & ItalianScript & \ttlprop{hm:ItalianScript} \\
\hline
ביזנטי / Byzantine & ByzantineScript & \ttlprop{hm:ByzantineScript} \\
\hline
תימני / Yemenite & YemeniteScript & \ttlprop{hm:YemeniteScript} \\
\hline
מזרחי / Oriental & OrientalScript & \ttlprop{hm:OrientalScript} \\
\hline
פרסי / Persian & PersianScript & \ttlprop{hm:PersianScript} \\
\hline
\end{longtable}

\subsection{מצבי כתב}

\begin{longtable}{|p{4cm}|L{5cm}|L{7cm}|}
\hline
\textbf{מונחים} & \textbf{ערך TTL} & \textbf{URI} \\
\hline
\endfirsthead
\hline
\textbf{מונחים} & \textbf{ערך TTL} & \textbf{URI} \\
\hline
\endhead
\hline
\endfoot
מרובע / Square & SquareScript & \ttlprop{hm:SquareScript} \\
\hline
בינוני / Semi-cursive & SemicursiveScript & \ttlprop{hm:SemicursiveScript} \\
\hline
רהוט / Cursive & CursiveScript & \ttlprop{hm:CursiveScript} \\
\hline
\end{longtable}

\subsection{היררכיות קאנוניות}

\begin{longtable}{|p{4cm}|L{5cm}|L{7cm}|}
\hline
\textbf{מונחים} & \textbf{ערך TTL} & \textbf{URI} \\
\hline
\endfirsthead
\hline
\textbf{מונחים} & \textbf{ערך TTL} & \textbf{URI} \\
\hline
\endhead
\hline
\endfoot
תנ"ך / Bible & Bible\_hierarchy & \ttlprop{hm:Bible\_hierarchy} \\
\hline
משנה / Mishnah & Mishnah\_hierarchy & \ttlprop{hm:Mishnah\_hierarchy} \\
\hline
תלמוד בבלי & Talmud\_Bavli\_hierarchy & \ttlprop{hm:Talmud\_Bavli\_hierarchy} \\
\hline
תלמוד ירושלמי & Talmud\_Yerushalmi\_hierarchy & \ttlprop{hm:Talmud\_Yerushalmi\_hierarchy} \\
\hline
משנה תורה / רמב"ם & Mishneh\_Torah\_hierarchy & \ttlprop{hm:Mishneh\_Torah\_hierarchy} \\
\hline
שולחן ערוך & Shulchan\_Aruch\_hierarchy & \ttlprop{hm:Shulchan\_Aruch\_hierarchy} \\
\hline
זוהר & Zohar\_hierarchy & \ttlprop{hm:Zohar\_hierarchy} \\
\hline
מדרש & Midrash\_hierarchy & \ttlprop{hm:Midrash\_hierarchy} \\
\hline
\end{longtable}

\subsection{סוגי מצב פיסי}

\begin{longtable}{|p{4cm}|L{5cm}|L{7cm}|}
\hline
\textbf{מונחים} & \textbf{ערך TTL} & \textbf{עברית} \\
\hline
\endfirsthead
\hline
\textbf{מונחים} & \textbf{ערך TTL} & \textbf{עברית} \\
\hline
\endhead
\hline
\endfoot
excellent / מצוין & Excellent & מצוין \\
\hline
good / טוב & Good & טוב \\
\hline
fair / סביר & Fair & סביר \\
\hline
poor / גרוע & Poor & גרוע \\
\hline
fragmentary / חלקי & Fragmentary & חלקי/קטעי \\
\hline
\end{longtable}

\subsection{סוגי ניקוד}

\begin{longtable}{|p{4cm}|L{5cm}|L{7cm}|}
\hline
\textbf{מונחים} & \textbf{ערך TTL} & \textbf{URI} \\
\hline
\endfirsthead
\hline
\textbf{מונחים} & \textbf{ערך TTL} & \textbf{URI} \\
\hline
\endhead
\hline
\endfoot
טברייני / Tiberian & TiberianVocalization & \ttlprop{hm:TiberianVocalization} \\
\hline
בבלי / Babylonian & BabylonianVocalization & \ttlprop{hm:BabylonianVocalization} \\
\hline
ארץ-ישראלי / Palestinian & PalestinianVocalization & \ttlprop{hm:PalestinianVocalization} \\
\hline
ללא ניקוד / None & NoVocalization & \ttlprop{hm:NoVocalization} \\
\hline
\end{longtable}

\subsection{סוגי כריכה}

\begin{longtable}{|p{4cm}|L{5cm}|L{7cm}|}
\hline
\textbf{מונחים} & \textbf{ערך TTL} & \textbf{עברית} \\
\hline
\endfirsthead
\hline
\textbf{מונחים} & \textbf{ערך TTL} & \textbf{עברית} \\
\hline
\endhead
\hline
\endfoot
original & OriginalBinding & כריכה מקורית \\
\hline
rebinding & Rebinding & כריכה מחדש \\
\hline
conservation & ConservationBinding & כריכת שימור \\
\hline
unbound & Unbound & ללא כריכה \\
\hline
\end{longtable}

\subsection{סוגי פורמט תאריך )v1.4(}

\begin{longtable}{|p{4cm}|L{5cm}|L{7cm}|}
\hline
\textbf{מונחים} & \textbf{ערך TTL} & \textbf{דוגמאות} \\
\hline
\endfirsthead
\hline
\textbf{מונחים} & \textbf{ערך TTL} & \textbf{דוגמאות} \\
\hline
\endhead
\hline
\endfoot
gregorian\_year / שנה גרגוריאנית & GregorianYear & 1407, 1523, 1890 \\
\hline
hebrew\_year / שנה עברית & HebrewYear & תשפ״ד, ה'קס״ז, שנת תש״ה \\
\hline
full\_date / תאריך מלא & FullDate & 15/03/1407, כ״ג אדר תשפ״ד \\
\hline
unstructured / טקסט חופשי & UnstructuredDate & "המאה ה-15", "ראשית המאה" \\
\hline
\end{longtable}

\subsection{סוגי קישוט}

\begin{longtable}{|p{4cm}|L{5cm}|L{7cm}|}
\hline
\textbf{מונחים} & \textbf{ערך TTL} & \textbf{עברית} \\
\hline
\endfirsthead
\hline
\textbf{מונחים} & \textbf{ערך TTL} & \textbf{עברית} \\
\hline
\endhead
\hline
\endfoot
full page illumination & FullPageIllumination & עיטור דף שלם \\
\hline
initial letters & InitialLetters & אותיות פותחות מעוטרות \\
\hline
marginal decoration & MarginalDecoration & עיטור שוליים \\
\hline
carpet page & CarpetPage & דף שטיח \\
\hline
micrography & Micrography & מיקרוגרפיה \\
\hline
\end{longtable}

\subsection{סוגי ז'אנר}

\begin{longtable}{|p{4cm}|L{5cm}|L{7cm}|}
\hline
\textbf{מונחים} & \textbf{ערך TTL} & \textbf{עברית} \\
\hline
\endfirsthead
\hline
\textbf{מונחים} & \textbf{ערך TTL} & \textbf{עברית} \\
\hline
\endhead
\hline
\endfoot
bible / תנ"ך & BiblicalText & טקסט תנ"כי \\
\hline
mishnah / משנה & MishnaicText & טקסט משנאי \\
\hline
talmud / תלמוד & TalmudicText & טקסט תלמודי \\
\hline
halakha / הלכה & HalachicText & טקסט הלכתי \\
\hline
kabbalah / קבלה & KabbalisticText & טקסט קבלי \\
\hline
philosophy / פילוסופיה & PhilosophicalText & טקסט פילוסופי \\
\hline
poetry / שירה & PoeticText & טקסט שירי \\
\hline
liturgy / תפילה & LiturgicalText & טקסט ליטורגי \\
\hline
medicine / רפואה & MedicalText & טקסט רפואי \\
\hline
astronomy / אסטרונומיה & AstronomicalText & טקסט אסטרונומי \\
\hline
grammar / דקדוק & GrammaticalText & טקסט דקדוקי \\
\hline
commentary / פירוש & CommentaryText & פירוש \\
\hline
\end{longtable}

\subsection{סוגי טקסט )מבני(}

סיווג מבני של טקסטים לפי מיקומם ותפקידם בכתב היד:

\begin{longtable}{|p{4cm}|L{5cm}|L{7cm}|}
\hline
\textbf{מונחים} & \textbf{ערך TTL} & \textbf{עברית} \\
\hline
\endfirsthead
\hline
\textbf{מונחים} & \textbf{ערך TTL} & \textbf{עברית} \\
\hline
\endhead
\hline
\endfoot
main\_text / טקסט עיקרי & MainText & טקסט עיקרי \\
\hline
commentary / פירוש & CommentaryText & טקסט פרשני \\
\hline
supercommentary / פירוש על פירוש & SupercommentaryText & פירוש על פירוש \\
\hline
gloss / ביאור & GlossText & ביאור קצר \\
\hline
responsum / תשובה & ResponsumText & תשובה הלכתית \\
\hline
additional / נוסף & AdditionalText & טקסט נוסף \\
\hline
marginal / שוליים & MarginalText & טקסט בשוליים \\
\hline
\end{longtable}

\subsection{שיטות פרשנות )v1.4(}

\begin{longtable}{|p{4cm}|L{5cm}|L{7cm}|}
\hline
\textbf{מונחים} & \textbf{ערך TTL} & \textbf{עברית} \\
\hline
\endfirsthead
\hline
\textbf{מונחים} & \textbf{ערך TTL} & \textbf{עברית} \\
\hline
\endhead
\hline
\endfoot
paleographic & PaleographicAnalysis & ניתוח פליאוגרפי \\
\hline
codicological & CodicologicalAnalysis & ניתוח קודיקולוגי \\
\hline
linguistic & LinguisticAnalysis & ניתוח לשוני \\
\hline
historical & HistoricalContextAnalysis & ניתוח הקשר היסטורי \\
\hline
comparative & ComparativeTextualAnalysis & ניתוח טקסטואלי השוואתי \\
\hline
ai & AIBasedAnalysis & ניתוח מבוסס AI \\
\hline
radiocarbon & RadiocarbonDating & תיארוך פחמן-14 \\
\hline
watermark & WatermarkAnalysis & ניתוח סימני מים \\
\hline
\end{longtable}

\subsection{קטגוריות נתונים )v1.4(}

\begin{longtable}{|p{4cm}|L{5cm}|L{7cm}|}
\hline
\textbf{קטגוריה} & \textbf{ערך TTL} & \textbf{עובדתי?} \\
\hline
\endfirsthead
\hline
\textbf{קטגוריה} & \textbf{ערך TTL} & \textbf{עובדתי?} \\
\hline
\endhead
\hline
\endfoot
מדידה פיזית & PhysicalMeasurement & כן \\
\hline
תוכן קולופון & ColophonContent & כן \\
\hline
תצפית חומר & MaterialObservation & כן \\
\hline
ייחוס תאריך & DateAttribution & לא \\
\hline
ייחוס מחברות & AuthorshipAttribution & לא \\
\hline
זיהוי סופר & ScribeIdentification & לא \\
\hline
ייחוס מקום & PlaceAttribution & לא \\
\hline
זיהוי יצירה & WorkIdentification & לא \\
\hline
יחס טקסטואלי & TextualRelationship & לא \\
\hline
\end{longtable}

\subsection{סוגי התערבויות סופר}

\begin{longtable}{|p{4cm}|L{5cm}|L{7cm}|}
\hline
\textbf{מונחים} & \textbf{ערך TTL} & \textbf{עברית} \\
\hline
\endfirsthead
\hline
\textbf{מונחים} & \textbf{ערך TTL} & \textbf{עברית} \\
\hline
\endhead
\hline
\endfoot
correction / תיקון & Correction\_type & תיקון \\
\hline
erasure / מחיקה & Erasure\_type & מחיקה \\
\hline
interlinear\_addition & Interlinear\_addition\_type & הוספה בין השורות \\
\hline
marginal\_gloss & Marginal\_gloss\_type & ביאור בשוליים \\
\hline
later\_hand / יד מאוחרת & Later\_hand\_type & יד מאוחרת \\
\hline
\end{longtable}

\subsection{סוגי חשיבות גרסה}

\begin{longtable}{|p{4cm}|L{5cm}|L{7cm}|}
\hline
\textbf{מונחים} & \textbf{ערך TTL} & \textbf{עברית} \\
\hline
\endfirsthead
\hline
\textbf{מונחים} & \textbf{ערך TTL} & \textbf{עברית} \\
\hline
\endhead
\hline
\endfoot
orthographic / גרסה כתיבית & Orthographic\_variant & גרסה כתיבית \\
\hline
lexical / גרסה לקסיקלית & Lexical\_variant & גרסה לקסיקלית \\
\hline
syntactic / גרסה תחבירית & Syntactic\_variant & גרסה תחבירית \\
\hline
semantic / גרסה משמעותית & Semantic\_variant & גרסה משמעותית \\
\hline
addition / הוספה & Addition\_variant & הוספה \\
\hline
omission / השמטה & Omission\_variant & השמטה \\
\hline
\end{longtable}

\subsection{סטטוס יחידה}

\begin{longtable}{|p{4cm}|L{5cm}|L{7cm}|}
\hline
\textbf{מונחים} & \textbf{ערך TTL} & \textbf{עברית} \\
\hline
\endfirsthead
\hline
\textbf{מונחים} & \textbf{ערך TTL} & \textbf{עברית} \\
\hline
\endhead
\hline
\endfoot
core & CoreUnit\_status & יחידת ליבה \\
\hline
later\_addition & LaterAddition\_status & תוספת מאוחרת \\
\hline
binder\_addition & BinderAddition\_status & תוספת כורך \\
\hline
unrelated\_fragment & UnrelatedFragment\_status & שבר לא קשור \\
\hline
protective\_leaf & ProtectiveLeaf\_status & דף מגן \\
\hline
\end{longtable}

\subsection{סוגי היררכיה )v1.4(}

\begin{longtable}{|p{4cm}|L{5cm}|L{7cm}|}
\hline
\textbf{מונחים} & \textbf{ערך TTL} & \textbf{עברית} \\
\hline
\endfirsthead
\hline
\textbf{מונחים} & \textbf{ערך TTL} & \textbf{עברית} \\
\hline
\endhead
\hline
\endfoot
simple & SimpleHierarchy & היררכיה פשוטה \\
\hline
nested & NestedHierarchy & היררכיה מקוננת \\
\hline
complex & ComplexHierarchy & היררכיה מורכבת \\
\hline
fragmentary & FragmentaryHierarchy & היררכיה קטעית \\
\hline
\end{longtable}

\subsection{סוגי פרדיגמה )v1.4(}

\begin{longtable}{|p{4cm}|L{5cm}|L{7cm}|}
\hline
\textbf{מונחים} & \textbf{ערך TTL} & \textbf{עברית} \\
\hline
\endfirsthead
\hline
\textbf{מונחים} & \textbf{ערך TTL} & \textbf{עברית} \\
\hline
\endhead
\hline
\endfoot
bibliographic & BibliographicParadigm & פרדיגמה ביבליוגרפית \\
\hline
philological & PhilologicalParadigm & פרדיגמה פילולוגית \\
\hline
hybrid & HybridParadigm & פרדיגמה היברידית \\
\hline
\end{longtable}

\newpage
% =============================================================================
\section{דוגמה מלאה}
% =============================================================================

\subsection{רשומת MARC לדוגמה}

\begin{LTR}
\begin{verbatim}
001 990012345680205171
008 850101s1407    xx            000 0 heb d
041 0# $a heb $a arc
100 1# $a רמב"ם $d 1138-1204
245 10 $a משנה תורה $b ספר המדע
260 ## $a ספרד $c [1407]
300 ## $a 248 דפים $c 280 x 200 מ"מ
340 ## $a קלף
500 ## $a כתב ספרדי בינוני. יש תיקונים בשוליים.
500 ## $a קולופון: נשלם בשנת קס"ז לפ"ק
505 0# $a הלכות יסודי התורה )דפים 1-50( -- הלכות דעות )51-100(
510 4# $a Institute of Hebrew Manuscripts
561 ## $a אוסף ששון
563 ## $a כריכת עור מקורית
600 10 $a רמב"ם $d 1138-1204
650 #0 $a הלכה
655 #7 $a כתבי יד
700 1# $a יצחק בן משה $e סופר
856 40 $u https://www.nli.org.il/...
\end{verbatim}
\end{LTR}

\subsection{פלט TTL}

\begin{LTR}
\begin{verbatim}
@prefix hm: <http://www.ontology.org.il/HebrewManuscripts/2025-12-06#> .
@prefix lrmoo: <http://iflastandards.info/ns/lrm/lrmoo/> .
@prefix cidoc-crm: <http://www.cidoc-crm.org/cidoc-crm/> .
@prefix rdf: <http://www.w3.org/1999/02/22-rdf-syntax-ns#> .
@prefix rdfs: <http://www.w3.org/2000/01/rdf-schema#> .
@prefix xsd: <http://www.w3.org/2001/XMLSchema#> .

# Manuscript
hm:MS_990012345680205171 a lrmoo:F4_Manifestation_Singleton,
                           hm:Codicological_Unit ;
  rdfs:label "משנה תורה"@he ;
  hm:external_identifier_nli "990012345680205171"^^xsd:string ;
  hm:has_number_of_folios 248 ;
  hm:has_height_mm 280 ;
  hm:has_width_mm 200 ;
  hm:has_material hm:Parchment ;
  hm:has_script_type hm:SepharadicScript ;
  hm:has_script_mode hm:SemicursiveScript ;
  hm:ownership_history "אוסף ששון"@he ;
  hm:has_digital_representation_url "https://www.nli.org.il/..."^^xsd:anyURI ;
  lrmoo:R4_embodies hm:Expression_משנה_תורה_in_990012345680205171 ;
  hm:has_colophon hm:Colophon_990012345680205171 ;
  hm:has_binding hm:Binding_990012345680205171 ;
  hm:has_cataloging_perspective hm:CatalogingView_990012345680205171 ;
  hm:has_evidence_chain hm:EvidenceChain_990012345680205171 .

# Work
hm:Work_משנה_תורה_by_רמבם a lrmoo:F1_Work ;
  rdfs:label "משנה תורה"@he ;
  hm:has_title "משנה תורה"@he ;
  hm:has_title "משנה תורה : ספר המדע"@he ;
  cidoc-crm:P2_has_type hm:Subject_כתבי_יד ;
  cidoc-crm:P129_is_about hm:Person_רמבם, hm:Subject_הלכה .

# Expression
hm:Expression_משנה_תורה_in_990012345680205171 a lrmoo:F2_Expression ;
  rdfs:label "משנה תורה (in MS 990012345680205171)"@he ;
  lrmoo:R3_is_realised_in hm:Work_משנה_תורה_by_רמבם ;
  cidoc-crm:P72_has_language hm:Hebrew, hm:Aramaic .

# Author
hm:Person_רמבם a cidoc-crm:E21_Person ;
  rdfs:label "רמב\"ם"@he ;
  cidoc-crm:P82a_begin_of_the_begin 1138 ;
  cidoc-crm:P82b_end_of_the_end 1204 .

# Work Creation
hm:WorkCreation_משנה_תורה_by_רמבם a lrmoo:F27_Work_Creation ;
  lrmoo:R16_created hm:Work_משנה_תורה_by_רמבם ;
  cidoc-crm:P14_carried_out_by hm:Person_רמבם .

# Scribe
hm:Person_יצחק_בן_משה a cidoc-crm:E21_Person ;
  rdfs:label "יצחק בן משה"@he .

# Production
hm:Production_990012345680205171 a cidoc-crm:E12_Production ;
  lrmoo:R27_materialized hm:MS_990012345680205171 ;
  cidoc-crm:P7_took_place_at hm:Place_ספרד ;
  cidoc-crm:P4_has_time-span hm:TimeSpan_1407 ;
  cidoc-crm:P14_carried_out_by hm:Person_יצחק_בן_משה .

hm:Place_ספרד a cidoc-crm:E53_Place ;
  rdfs:label "ספרד"@he .

hm:TimeSpan_1407 a cidoc-crm:E52_Time-Span ;
  rdfs:label "1407" ;
  hm:earliest_possible_date 1407 ;
  hm:latest_possible_date 1407 .

# Colophon
hm:Colophon_990012345680205171 a hm:Colophon ;
  hm:colophon_text "קולופון: נשלם בשנת קס\"ז לפ\"ק"@he .

# Binding
hm:Binding_990012345680205171 a hm:Binding ;
  rdfs:comment "כריכת עור מקורית"@he .

# Content Works
hm:Work_הלכות_יסודי_התורה a lrmoo:F1_Work ;
  hm:has_title "הלכות יסודי התורה"@he .

hm:Expression_הלכות_יסודי_התורה_in_990012345680205171 a lrmoo:F2_Expression ;
  lrmoo:R3_is_realised_in hm:Work_הלכות_יסודי_התורה ;
  hm:has_folio_range "1-50" .

hm:MS_990012345680205171 lrmoo:R4_embodies 
  hm:Expression_הלכות_יסודי_התורה_in_990012345680205171 .

# Cataloging View (v1.4)
hm:CatalogingView_990012345680205171 a hm:CatalogingView ;
  rdfs:label "Cataloging view for MS 990012345680205171" ;
  hm:cataloging_work hm:Work_משנה_תורה_by_רמבם ;
  hm:cataloging_expression hm:Expression_משנה_תורה_in_990012345680205171 ;
  hm:is_primary_paradigm true ;
  hm:paradigm_note "Data derived from NLI catalog using 
                    bibliographic Work-Expression model"@en .

# Evidence Chain (v1.4)
hm:EvidenceChain_990012345680205171 a hm:EvidenceChain ;
  rdfs:label "Evidence chain for MS 990012345680205171" ;
  hm:evidence_step hm:CatalogStep_990012345680205171 .

hm:CatalogStep_990012345680205171 a hm:EvidenceStep ;
  hm:reasoning_text "Data imported from National Library 
                     of Israel MARC catalog"@en .

hm:MS_990012345680205171 
  hm:attribution_source hm:CatalogAttribution ;
  hm:has_epistemological_status hm:CatalogInherited ;
  hm:is_factual false .

# Languages
hm:Hebrew a cidoc-crm:E56_Language ;
  rdfs:label "Hebrew" .

hm:Aramaic a cidoc-crm:E56_Language ;
  rdfs:label "Aramaic" .

# Materials
hm:Parchment a cidoc-crm:E57_Material ;
  rdfs:label "Parchment"@en .

# Catalog Reference
hm:Catalog_Institute_of_Hebrew_Manuscripts a lrmoo:F3_Manifestation ;
  rdfs:label "Institute of Hebrew Manuscripts" .

hm:MS_990012345680205171 cidoc-crm:P70i_is_documented_in 
  hm:Catalog_Institute_of_Hebrew_Manuscripts .
\end{verbatim}
\end{LTR}

\newpage
\begin{center}
\textbf{נספח: תיעוד האונטולוגיה}\\
התיעוד להלן נלקח מתוך \textenglish{hebrew\_ontology\_documentation.tex}
\end{center}

\section{מבוא}
%====================================================================

אונטולוגיה זו פותחה לתיאור כתבי יד עבריים, והיא מבוססת על תקנים בינלאומיים מוכרים:

\begin{itemize}
\item \textbf{CIDOC-CRM} — מודל ייחוס מושגי למוזיאונים \textenglish{(ISO 21127)}
\item \textbf{FRBRoo/LRMoo} — מודל ספרני מונחה עצמים לתיאור ביבליוגרפי
\item \textbf{SKOS} — מערכת ארגון ידע פשוטה לאוצרות מילים
\item הרחבות מותאמות לכתבי יד עבריים
\end{itemize}

\textbf{הגישה העיקרית היא אירועית} — תיעוד מה קרה לכתב היד לאורך הזמן )יצירה, העברות בעלות, תיעוד(, מהיצירה המופשטת \textenglish{(Work)} לביטוי הטקסטואלי \textenglish{(Expression)} למופע הפיזי \textenglish{(Manifestation/Item)}.

\subsection{היררכיית FRBRoo/LRMoo הבסיסית}

\begin{verbatim}
F1_Work (יצירה מופשטת)
    ↓ R3_is_realised_in
F2_Expression (ביטוי טקסטואלי)
    ↓ R4i_is_embodied_in
F4_Manifestation_Singleton (עותק פיזי יחידאי)
    ≡ F5_Item (תואם LRMoo)
\end{verbatim}

\subsection{מבנה תלת-שכבתי \textenglish{(BU-CU-PU)}}

המבנה הפיזי של כתבי יד מיוצג בשלוש רמות:

\begin{itemize}
\item \textbf{BU} \textenglish{(Bibliographic Unit)} — יחידה ביבליוגרפית: הכרך הפיזי השלם
\item \textbf{CU} \textenglish{(Codicological Unit)} — יחידה קודיקולוגית: חיבור מובחן שנוצר בזמן/מקום מסוים
\item \textbf{PU} \textenglish{(Paleographical Unit)} — יחידה פלאוגרפית: חלק כתוב ביד כותב ספציפי
\end{itemize}

כל הרמות הן תת-מחלקות של \textenglish{F4\_Manifestation\_Singleton}, מקושרות ביחסי הרכבה \textenglish{R5\_has\_component}.

\subsection{מרחבי שמות (Namespaces)}

\begin{table}[h]
\centering
\begin{tabular}{|r|l|}
\hline
\textbf{קידומת} & \textbf{URI} \\
\hline
hm: & http://www.ontology.org.il/HebrewManuscripts/2025-12-06\# \\
lrmoo: & http://iflastandards.info/ns/lrm/lrmoo/ \\
cidoc-crm: & http://www.cidoc-crm.org/cidoc-crm/ \\
skos: & http://www.w3.org/2004/02/skos/core\# \\
\textcolor{newaddition}{nli:} & \textcolor{newaddition}{https://www.nli.org.il/en/authorities/} \\
\hline
\end{tabular}
\caption{מרחבי שמות בשימוש}
\end{table}

\subsection{\textcolor{newaddition}{אינטגרציה עם מאגר הרשויות הלאומי — מזל}}

\textcolor{newaddition}{החל מגרסה 1.5, המערכת משתלבת עם מאגר הרשויות הישראלי מזל \textenglish{(Mazal)} של הספרייה הלאומית.}

\textcolor{newaddition}{מזל הוא קובץ הרשויות הלאומי הישראלי המכיל כ-4.5 מיליון רשומות רשות, המספקות זיהוי ייחודי ואחיד עבור:}
\begin{itemize}
\item \textcolor{newaddition}{אישים — שמות אישיים}
\item \textcolor{newaddition}{מקומות — שמות גיאוגרפיים}
\item \textcolor{newaddition}{גופים — שמות תאגידים וארגונים}
\item \textcolor{newaddition}{יצירות — כותרות אחידות}
\end{itemize}

\textcolor{newaddition}{יתרונות האינטגרציה:}
\begin{itemize}
\item \textcolor{newaddition}{זיהוי חד-משמעי של ישויות — במקום URI מקומי, משתמשים ב-NLI ID רשמי}
\item \textcolor{newaddition}{קישוריות בינלאומית — מזל מקושר ל-VIAF, Wikidata, ORCID}
\item \textcolor{newaddition}{תמיכה רב-לשונית — עברית, לטינית, ערבית, קירילית}
\item \textcolor{newaddition}{עקביות קטלוגית — שימוש באותה צורת שם כמו בקטלוגים אחרים}
\end{itemize}

\textcolor{newaddition}{דוגמה — לפני ואחרי האינטגרציה:}

\begin{verbatim}
# לפני (URI מקומי):
hm:Person_רמבם a cidoc-crm:E21_Person ;
    rdfs:label "רמב\"ם"@he .

# אחרי (URI מזל):
<https://www.nli.org.il/en/authorities/987007388484005171>
    a cidoc-crm:E21_Person ;
    rdfs:label "רמב\"ם"@he ;
    hm:external_identifier_nli "987007388484005171" .
\end{verbatim}

\newpage

%====================================================================
\section{מחלקות (Classes)}
%====================================================================

%--------------------------------------------------------------------
\subsection{מחלקות LRMoo — מודל ביבליוגרפי}
%--------------------------------------------------------------------

\begin{longtable}{|r|p{10cm}|}
\hline
\textbf{מחלקה} & \textbf{תיאור} \\
\hline
\endhead

lrmoo:F1\_Work & 
יצירה מופשטת — הרעיון האינטלקטואלי של חיבור, ללא תלות בצורה טקסטואלית ספציפית.\\
\hline

lrmoo:F2\_Expression & 
ביטוי — מימוש ספציפי של יצירה: גרסה מסוימת עם ניסוח, שפה וצורה ייחודיים.\\
\hline

lrmoo:F3\_Manifestation & 
גילום — התגלמות פיזית של ביטוי; תבנית שממנה נעשים עותקים.\\
\hline

lrmoo:F4\_Manifestation\_Singleton & 
גילום יחידאי — אובייקט פיזי ייחודי, כלומר כתב היד הבודד.\\
\hline

lrmoo:F5\_Item & 
פריט — עותק פיזי ספציפי. \textbf{הערה:} \textenglish{F4\_Manifestation\_Singleton ≡ F5\_Item} — שקילות לתאימות בין \textenglish{FRBRoo} ישן ל-\textenglish{LRMoo} חדש.\\
\hline

lrmoo:F27\_Work\_Creation & 
אירוע יצירת יצירה — מתי ובידי מי נוצר החיבור המקורי.\\
\hline

lrmoo:F28\_Expression\_Creation & 
אירוע יצירת ביטוי — יצירת גרסה טקסטואלית ספציפית.\\
\hline

lrmoo:F30\_Manifestation\_Creation & 
אירוע יצירת גילום.\\
\hline

lrmoo:F32\_Item\_Production\_Event & 
אירוע ייצור פריט — העתקת כתב היד.\\
\hline
\end{longtable}

%--------------------------------------------------------------------
\subsection{מחלקות CIDOC-CRM — מורשת תרבותית}
%--------------------------------------------------------------------

\begin{longtable}{|r|p{10cm}|}
\hline
\textbf{מחלקה} & \textbf{תיאור} \\
\hline
\endhead

cidoc-crm:E1\_CRM\_Entity & 
ישות CRM — מחלקת בסיס לכל הישויות במודל CIDOC-CRM.\\
\hline

cidoc-crm:E18\_Physical\_Thing & 
דבר פיזי — כל חפץ פיזי, כולל כתבי יד ופריטים חומריים.\\
\hline

cidoc-crm:E39\_Actor & 
שחקן — מחלקת על לסוכנים (אנשים וקבוצות).\\
\hline

cidoc-crm:E35\_Title & 
כותרת — ישות כותרת נפרדת במודל.\\
\hline

cidoc-crm:E73\_Information\_Object & 
אובייקט מידע — מסמך/אובייקט מידע המשמש לתיעוד ידע.\\
\hline

cidoc-crm:E7\_Activity & 
פעילות — מחלקת בסיס לאירועים ופעולות.\\
\hline

cidoc-crm:E8\_Acquisition & 
רכישה — אירוע רכישת בעלות על חפץ.\\
\hline

cidoc-crm:E10\_Transfer\_of\_Custody & 
העברת משמורת — העברת אחזקה פיזית של כתב יד.\\
\hline

cidoc-crm:E12\_Production & 
הפקה — יצירת חפץ פיזי )כולל העתקת כתב יד(.\\
\hline

cidoc-crm:E21\_Person & 
אדם — אדם בודד (מחבר, סופר, בעלים).\\
\hline

cidoc-crm:E22\_Human-Made\_Object & 
אובייקט מעשה אדם — חפצים פיזיים שנוצרו בידי אדם.\\
\hline

cidoc-crm:E33\_Linguistic\_Object & 
אובייקט לשוני — טקסטים וביטויים לשוניים.\\
\hline

cidoc-crm:E42\_Identifier & 
מזהה — סימני מדף, מספרי קטלוג.\\
\hline

cidoc-crm:E52\_Time-Span & 
טווח זמן — תקופות ותאריכים.\\
\hline

cidoc-crm:E53\_Place & 
מקום — מיקומים גיאוגרפיים.\\
\hline

cidoc-crm:E55\_Type & 
סוג — קטגוריות וסיווגים.\\
\hline

cidoc-crm:E56\_Language & 
שפה — שפות )עברית, ארמית, ערבית-יהודית(.\\
\hline

cidoc-crm:E57\_Material & 
חומר — חומרי כתיבה )קלף, נייר(.\\
\hline

cidoc-crm:E74\_Group & 
קבוצה — קבוצות אנשים (מוסדות, משפחות).\\
\hline

cidoc-crm:E78\_Curated\_Holding & 
אוסף — אוספים מתוחזקים.\\
\hline

cidoc-crm:E28\_Conceptual\_Object & 
אובייקט מושגי — מושגים מופשטים.\\
\hline
\end{longtable}

%--------------------------------------------------------------------
\subsection{מחלקות כלליות (OWL/RDF)}
%--------------------------------------------------------------------

\begin{longtable}{|r|p{10cm}|}
\hline
\textbf{מחלקה} & \textbf{תיאור} \\
\hline
\endhead

owl:Thing & ישות כללית — טיפוס על לכל ישות RDF/OWL.\\
\hline
\end{longtable}

%--------------------------------------------------------------------
\subsection{מחלקות ליבה לכתבי יד עבריים}
%--------------------------------------------------------------------

\begin{longtable}{|r|p{10cm}|}
\hline
\textbf{מחלקה} & \textbf{תיאור} \\
\hline
\endhead

hm:Bibliographic\_Unit & 
\textbf{יחידה ביבליוגרפית} — רמת תיאור הקטלוג; מתאים לסימן מדף אחד.\\
\hline

hm:Codicological\_Unit & 
\textbf{יחידה קודיקולוגית} — חלק שיוצר כיחידה פיזית; קוויונים עם מאפייני חומר משותפים.\\
\hline

hm:Paleographical\_Unit & 
\textbf{יחידה פליאוגרפית} — קטע עם מאפייני כתב עקביים; עבודת יד אחת.\\
\hline

hm:Colophon & 
\textbf{קולופון} — הצהרת הסופר בסוף הטקסט עם מידע על ההעתקה.\\
\hline

hm:HebrewScriptType & 
\textbf{סוג כתב עברי} — מחלקת בסיס לסיווגי כתב.\\
\hline

hm:TypeScriptType & 
\textbf{סוג כתב אזורי} — אשכנזי, ספרדי, ביזנטי, תימני וכו'.\\
\hline

hm:ModeScriptType & 
\textbf{מצב כתב} — מרובע, בינוני, רהוט.\\
\hline

hm:Binding & 
\textbf{כריכה} — כריכת כתב היד כאובייקט נפרד.\\
\hline

hm:Decoration & 
\textbf{עיטור} — אלמנטים דקורטיביים.\\
\hline

hm:Marginalia & 
\textbf{הערות שוליים} — הערות מאוחרות בשולי הדף.\\
\hline

hm:Watermark & 
\textbf{סימן מים} — סימני מים בנייר לצורך תיארוך ומיקום.\\
\hline

hm:DecorationType & 
\textbf{סוג עיטור} — סוגי קישוט (איור דף מלא, אותיות פותחות, מיקרוגרפיה).\\
\hline

hm:ConditionType & 
\textbf{סוג מצב} — מצב פיזי (מצוין, טוב, בינוני, גרוע, קטעי).\\
\hline

hm:VocalizationType & 
\textbf{סוג ניקוד} — טברייני, בבלי, ארץ-ישראלי, ללא.\\
\hline

hm:BindingType & 
\textbf{סוג כריכה} — מקורית, כריכה מחדש, כריכת שימור.\\
\hline

\new{hm:DateFormatType} & 
\new{\textbf{סוג פורמט תאריך} — שנה גרגוריאנית, שנה עברית, תאריך מלא, טקסט חופשי.}\\
\hline

hm:ParticipationRole & 
\textbf{תפקיד השתתפות} — תפקיד אנשים (סופר, מחבר, מאייר, מתרגם, בעלים).\\
\hline

hm:SubjectType & 
\textbf{נושא} — סיווג נושאי (מקרא, תלמוד, הלכה, קבלה, פילוסופיה).\\
\hline

hm:Reference\_Event & 
\textbf{אירוע הפניה} — אירוע קטלוג או הפניה במחקר.\\
\hline
\end{longtable}

%--------------------------------------------------------------------
\subsection{\new{ ודאות וייחוס}}
%--------------------------------------------------------------------

\begin{longtable}{|r|p{10cm}|}
\hline
\textbf{מחלקה} & \textbf{תיאור} \\
\hline
\endhead

\newclass{hm:CertaintyLevel} & 
\new{רמת ודאות — דרגת ביטחון בייחוסים (ודאי, סביר, אפשרי, לא ודאי).}\\
\hline

\newclass{hm:AttributionSource} & 
\new{מקור ייחוס — מומחה, AI, קטלוג, קולופון, פליאוגרפי.}\\
\hline

\newclass{hm:MultiVolumeSet} & 
\new{סט רב-כרכי — קבוצת כרכים השייכים יחד.}\\
\hline

\newclass{hm:AnthologyStructure} & 
\new{מבנה אנתולוגיה — כתב יד עם סידור אנתולוגי מכוון.}\\
\hline

\newclass{hm:ScribalIntervention} & 
\new{התערבות סופר — תיקונים, הוספות, שינויים.}\\
\hline

\newclass{hm:LostManuscript} & 
\new{כתב יד אבוד — כתב יד שאינו קיים אך ידוע מעותקים.}\\
\hline

\newclass{hm:HypotheticalExemplar} & 
\new{דוגמה היפותטית — מקור משוחזר על בסיס ניתוח טקסטואלי.}\\
\hline
\end{longtable}

%--------------------------------------------------------------------
\subsection{\new{ מסורת טקסט וגרנולריות}}
%--------------------------------------------------------------------

\begin{longtable}{|r|p{10cm}|}
\hline
\textbf{מחלקה} & \textbf{תיאור} \\
\hline
\endhead

\newclass{hm:TextTradition} & 
\new{מסורת טקסט — זרם העברה של טקסט, ללא תלות במושג "יצירה מקורית".}\\
\hline

\newclass{hm:TransmissionWitness} & 
\new{עד העברה — כתב יד כעד למסורת טקסטואלית.}\\
\hline

\newclass{hm:TextualVariant} & 
\new{וריאנט טקסטואלי — קריאה חלופית ספציפית.}\\
\hline

\newclass{hm:TextLocation} & 
\new{מיקום בטקסט — מיקום ספציפי בכתב היד.}\\
\hline

\newclass{hm:FolioLocation} & 
\new{מיקום דף — רמת דף.}\\
\hline

\newclass{hm:LineLocation} & 
\new{מיקום שורה — רמת שורה.}\\
\hline

\newclass{hm:WordLocation} & 
\new{מיקום מילה — רמת מילה.}\\
\hline

\newclass{hm:CharacterLocation} & 
\new{מיקום תו — רמת תו לניתוח פליאוגרפי.}\\
\hline

\newclass{hm:CanonicalReference} & 
\new{הפניה קנונית — הפניה למבנה טקסט קנוני.}\\
\hline

\newclass{hm:BiblicalReference} & 
\new{הפניה מקראית — ספר/פרק/פסוק.}\\
\hline

\newclass{hm:MishnaicReference} & 
\new{הפניה משנאית — מסכת/פרק/משנה.}\\
\hline

\newclass{hm:TalmudicReference} & 
\new{הפניה תלמודית — מסכת/דף/עמוד.}\\
\hline

\newclass{hm:HalachicReference} & 
\new{הפניה הלכתית — למשנה תורה או שולחן ערוך.}\\
\hline

\newclass{hm:CanonicalHierarchyType} & 
\new{סוג היררכיה קנונית — סיווג מבני טקסט (מקרא, משנה וכו').}\\
\hline

\newclass{hm:UnitStatusType} & 
\new{סטטוס יחידה — האם היחידה חלק מהליבה או תוספת זרה.}\\
\hline

\newclass{hm:TextSpan} & 
\new{טווח טקסט — טווח ממיקום א' למיקום ב'.}\\
\hline

\newclass{hm:IIIFRegion} & 
\new{אזור IIIF — הפניה לאזור תמונה בתקן IIIF.}\\
\hline

\newclass{hm:StemmaPosition} & 
\new{מיקום בסטמה — מיקום כתב יד בעץ סטמטי.}\\
\hline

\newclass{hm:StemmaHypothesis} & 
\new{השערת סטמה — השערה מדעית על יחסי טקסט.}\\
\hline

\newclass{hm:TextCorrection} & 
\new{תיקון טקסט — תיקון בטקסט על ידי סופר.}\\
\hline

\newclass{hm:MarginalAddition} & 
\new{הוספה בשוליים — טקסט שנוסף בשולי הדף.}\\
\hline

\newclass{hm:HandChange} & 
\new{חילופי יד — נקודה בה יד אחת מסתיימת ואחרת מתחילה.}\\
\hline

\newclass{hm:AnthologyPosition} & 
\new{מיקום באנתולוגיה — סדר ומיקום טקסט באנתולוגיה.}\\
\hline
\end{longtable}

%--------------------------------------------------------------------
\subsection{\new{פרדיגמה כפולה}}
%--------------------------------------------------------------------

אלה מחלקות חדשות שנוספו לפתרון המתח הפילוסופי בין הגישה הביבליוגרפית-מעשית לגישה הפילולוגית החדשה:

\begin{longtable}{|r|p{10cm}|}
\hline
\textbf{מחלקה} & \textbf{תיאור} \\
\hline
\endhead

\newclass{hm:ManuscriptView} & 
\new{\textbf{תצוגת כתב יד} — מחלקה מופשטת לתצוגות פרדיגמטיות שונות.}\\
\hline

\newclass{hm:CatalogingView} & 
\new{\textbf{תצוגת קטלוג} — תצוגה ביבליוגרפית-מעשית לפי מודל Work-Expression.}\\
\hline

\newclass{hm:PhilologicalView} & 
\new{\textbf{תצוגה פילולוגית} — תצוגה של פילולוגיה חדשה הרואה כל כתב יד כאירוע תרבותי ייחודי.}\\
\hline

\newclass{hm:ParadigmBridge} & 
\new{\textbf{גשר פרדיגמות} — קישור מפורש בין Work ל-TextTradition.}\\
\hline

\newclass{hm:ViewType} & 
\new{\textbf{סוג תצוגה} — סיווג סוגי תצוגה.}\\
\hline
\end{longtable}

%--------------------------------------------------------------------
\subsection{\new{  מבנה CU מקונן}}
%--------------------------------------------------------------------

מחלקות לתמיכה במבנים קודיקולוגיים מורכבים עם עומק שרירותי:

\begin{longtable}{|r|p{10cm}|}
\hline
\textbf{מחלקה} & \textbf{תיאור} \\
\hline
\endhead

\newclass{hm:CodicologicalHierarchy} & 
\new{\textbf{היררכיה קודיקולוגית} — מתאר את המבנה הכולל של יחידות מקוננות.}\\
\hline

\newclass{hm:CompositeCodicologicalUnit} & 
\new{\textbf{יחידה קודיקולוגית מורכבת} — יחידה המכילה תת-יחידות.}\\
\hline

\newclass{hm:AtomicCodicologicalUnit} & 
\new{\textbf{יחידה קודיקולוגית אטומית} — יחידה עלה ללא חלוקות משנה.}\\
\hline

\newclass{hm:HierarchyType} & 
\new{\textbf{סוג היררכיה} — סיווג מורכבות (פשוטה, מקוננת, מורכבת, פרגמנטרית).}\\
\hline
\end{longtable}

%--------------------------------------------------------------------
\subsection{\new{  מסגרת אפיסטמולוגית}}
%--------------------------------------------------------------------

מחלקות להבחנה בין עובדות לפרשנויות — מענה להערת פרופ' משה לביא:

\begin{longtable}{|r|p{10cm}|}
\hline
\textbf{מחלקה} & \textbf{תיאור} \\
\hline
\endhead

\newclass{hm:EpistemologicalStatus} & 
\new{\textbf{סטטוס אפיסטמולוגי} — סיווג נתונים לפי אופיים האפיסטמולוגי.}\\
\hline

\newclass{hm:EvidenceChain} & 
\new{\textbf{שרשרת ראיות} — שרשרת צעדי ראיות התומכים בטענה.}\\
\hline

\newclass{hm:EvidenceStep} & 
\new{\textbf{צעד ראיות} — צעד בודד בשרשרת ראיות.}\\
\hline

\newclass{hm:DirectObservationStep} & 
\new{\textbf{צעד תצפית ישירה} — ראיות מתצפית ישירה על כתב היד.}\\
\hline

\newclass{hm:InterpretationStep} & 
\new{\textbf{צעד פרשנות} — ראיות הכוללות פרשנות חוקר.}\\
\hline

\newclass{hm:ComputationalStep} & 
\new{\textbf{צעד חישובי} — ראיות שנגזרו מניתוח חישובי.}\\
\hline

\newclass{hm:DataCategory} & 
\new{\textbf{קטגוריית נתונים} — סיווג נתוני כתב יד לפי מקור.}\\
\hline

\newclass{hm:InterpretationMethod} & 
\new{\textbf{שיטת פרשנות} — שיטה מדעית לפרשנות.}\\
\hline
\end{longtable}

\newpage

%====================================================================
\section{תכונות אובייקט (Object Properties)}
%====================================================================

%--------------------------------------------------------------------
\subsection{תכונות LRMoo}
%--------------------------------------------------------------------

\begin{longtable}{|r|p{8cm}|}
\hline
\textbf{תכונה} & \textbf{תיאור} \\
\hline
\endhead

lrmoo:R3\_is\_realised\_in & מקשר Work ל-Expression שלו.\\
\hline
lrmoo:R4i\_is\_embodied\_in & מקשר Expression לגילום \textenglish{(F3/F4 Manifestation)} שבו הוא מגולם.\\
\hline
lrmoo:R4\_embodies & מקשר Manifestation ל-Expression שהוא מגלם.\\
\hline
lrmoo:R7\_exemplifies & עותק יחידאי מגלם מהדורה כללית.\\
\hline
lrmoo:R10\_has\_member & בת-יצירה כוללת יצירות (Work כולל Work).\\
\hline
lrmoo:R5\_has\_component & מקשר כתבי יד מורכבים לרכיביהם.\\
\hline
lrmoo:R5i\_is\_component\_of & יחס הפוך של R5\_has\_component )רכיב הוא חלק מכתב יד מורכב(.\\
\hline
lrmoo:R5i\_is\_component\_by & שם חלופי שמופיע בתיעוד המקורי ליחס ההופכי של R5\_has\_component.\\
\hline
lrmoo:R69\_has\_physical\_form & צורה פיזית )קישור לגנרליזציה של תבנית/צורה(.\\
\hline
lrmoo:R16\_created & מקשר אירוע יצירה ליצירה.\\
\hline
lrmoo:R17\_created & מקשר אירוע יצירת ביטוי ל-Expression.\\
\hline
lrmoo:R24\_created & מקשר אירוע יצירת גילום ל-Manifestation.\\
\hline
lrmoo:R27\_materialized & מקשר אירוע הפקה/העתקה לכתב יד יחידאי.\\
\hline
\end{longtable}

%--------------------------------------------------------------------
\subsection{תכונות CIDOC-CRM}
%--------------------------------------------------------------------

\begin{longtable}{|r|p{8cm}|}
\hline
\textbf{תכונה} & \textbf{תיאור} \\
\hline
\endhead

cidoc-crm:P2\_has\_type & שיוך לטיפוס/סוג.\\
\hline
cidoc-crm:P102\_has\_title & כותרת.\\
\hline
cidoc-crm:P129\_is\_about & נושא.\\
\hline
cidoc-crm:P1\_is\_identified\_by & מזהה.\\
\hline
cidoc-crm:P4\_has\_time-span & תקופת זמן.\\
\hline
cidoc-crm:P7\_took\_place\_at & מקום אירוע.\\
\hline
cidoc-crm:P14\_carried\_out\_by & מבצע פעילות.\\
\hline
cidoc-crm:P16\_used\_specific\_object & שימוש בחפץ/אובייקט ספציפי באירוע.\\
\hline
cidoc-crm:P23\_transferred\_title\_to & בעלים חדש באירוע רכישה.\\
\hline
cidoc-crm:P24\_transferred\_title\_from & בעלים קודם באירוע רכישה.\\
\hline
cidoc-crm:P28\_custody\_surrendered\_by & מסר משמורת באירוע העברת משמורת.\\
\hline
cidoc-crm:P29\_custody\_received\_by & קיבל משמורת באירוע העברת משמורת.\\
\hline
cidoc-crm:P70i\_is\_documented\_in & כתב יד מתועד בקטלוג/מסמך.\\
\hline
cidoc-crm:P72\_has\_language & שפה.\\
\hline
\end{longtable}

%--------------------------------------------------------------------
\subsection{תכונות ליבה}
%--------------------------------------------------------------------

\begin{longtable}{|r|p{8cm}|}
\hline
\textbf{תכונה} & \textbf{תיאור} \\
\hline
\endhead

hm:has\_script\_type & סוג כתב אזורי.\\
\hline
hm:has\_Script\_Mode & מצב כתב (מרובע, בינוני, רהוט).\\
\hline
hm:has\_colophon & מקשר לקולופון.\\
\hline
hm:has\_date\_of\_creation & תאריך יצירה.\\
\hline
\new{hm:has\_date\_format} & \new{סוג פורמט התאריך (GregorianYear, HebrewYear, FullDate, UnstructuredDate).}\\
\hline
hm:has\_material & חומר כתיבה.\\
\hline
hm:P44\_has\_condition & מצב פיזי/מצב שימור של כתב היד.\\
\hline
hm:has\_binding & מקשר לכריכה.\\
\hline
hm:has\_decoration & מקשר לעיטור.\\
\hline
hm:has\_marginalia & מקשר להערות שוליים.\\
\hline
hm:copied\_from & מקשר כתב יד לאב-טיפוס שלו.\\
\hline
hm:is\_translation\_of & יחס תרגום בין יצירות.\\
\hline
hm:has\_commentary\_on & יחס פירוש בין יצירות.\\
\hline

hm:has\_vocalization & סוג ניקוד (טברייני, בבלי, ארץ-ישראלי).\\
\hline

hm:has\_cantillation & טעמי המקרא.\\
\hline

hm:mentions\_date & מקשר קולופון לתאריך המוזכר בו.\\
\hline

hm:mentions\_place & מקשר קולופון למקום המוזכר בו.\\
\hline

hm:mentions\_scribe & מקשר קולופון לסופר המוזכר בו.\\
\hline

hm:mentions\_person & שם אדם המוזכר בכתב היד.\\
\hline
\end{longtable}

%--------------------------------------------------------------------
\subsection{\new{תכונות   — ודאות וייחוס}}
%--------------------------------------------------------------------

\begin{longtable}{|r|p{8cm}|}
\hline
\textbf{תכונה} & \textbf{תיאור} \\
\hline
\endhead

\newprop{hm:has\_certainty} & \new{רמת ודאות לכל ייחוס.}\\
\hline
\newprop{hm:attribution\_source} & \new{סוג מקור הייחוס.}\\
\hline
\newprop{hm:attributed\_by} & \new{מי ביצע את הייחוס.}\\
\hline
\newprop{hm:is\_variant\_of} & \new{יחס וריאנט בין יצירות.}\\
\hline
\newprop{hm:possibly\_realises} & \new{זיהוי Work לא ודאי.}\\
\hline
\newprop{hm:is\_abridgment\_of} & \new{גרסה מקוצרת.}\\
\hline
\newprop{hm:parallels} & \new{תוכן מקביל עם יחס לא ברור.}\\
\hline
\newprop{hm:has\_textual\_overlap\_with} & \new{חפיפה טקסטואלית (סימטרי).}\\
\hline
\newprop{hm:is\_volume\_of} & \new{מקשר כרך לסט שלו.}\\
\hline
\newprop{hm:has\_scribal\_intervention} & \new{מקשר להתערבות סופר.}\\
\hline
\newprop{hm:is\_copy\_of\_lost} & \new{מקשר לכתב יד אבוד.}\\
\hline

\newprop{hm:known\_through} & \new{כ"י אבוד ידוע דרך עותק קיים.}\\
\hline

\newprop{hm:is\_adaptation\_of} & \new{עיבוד של יצירה אחרת.}\\
\hline

\newprop{hm:duplicates\_part\_of} & \new{מכפיל חלק מביטוי אחר.}\\
\hline

\newprop{hm:has\_anthology\_position} & \new{מיקום טקסט באנתולוגיה.}\\
\hline

\newprop{hm:intervention\_location} & \new{מיקום התערבות בטקסט.}\\
\hline

\newprop{hm:intervention\_by} & \new{מבצע ההתערבות.}\\
\hline
\end{longtable}

%--------------------------------------------------------------------
\subsection{\new{תכונות   — מסורת טקסט}}
%--------------------------------------------------------------------

\begin{longtable}{|r|p{8cm}|}
\hline
\textbf{תכונה} & \textbf{תיאור} \\
\hline
\endhead

\newprop{hm:belongs\_to\_tradition} & \new{מקשר ביטוי למסורת טקסט.}\\
\hline
\newprop{hm:witnesses} & \new{כתב יד כעד למסורת.}\\
\hline
\newprop{hm:has\_variant\_reading} & \new{מקשר לקריאה חלופית.}\\
\hline
\newprop{hm:has\_location} & \new{קישור מיקום כללי.}\\
\hline
\newprop{hm:span\_start} & \new{תחילת טווח טקסט.}\\
\hline
\newprop{hm:span\_end} & \new{סוף טווח טקסט.}\\
\hline
\newprop{hm:covers\_canonical\_range} & \new{מקשר ביטוי לכיסוי קנוני.}\\
\hline
\newprop{hm:canonical\_hierarchy} & \new{סוג מערכת היררכיה.}\\
\hline
\newprop{hm:has\_unit\_status} & \new{סטטוס ליבה/זר.}\\
\hline

\newprop{hm:variant\_at\_location} & \new{גרסה במיקום ספציפי.}\\
\hline

\newprop{hm:tradition\_includes} & \new{מסורת כוללת ביטוי (הפוך של belongs\_to\_tradition).}\\
\hline

\newprop{hm:stemma\_parent} & \new{אב בסטמה טקסטואלית.}\\
\hline

\newprop{hm:stemma\_child} & \new{צאצא בסטמה.}\\
\hline

\newprop{hm:has\_stemma\_position} & \new{מיקום כתב יד בסטמה.}\\
\hline

\newprop{hm:iiif\_region} & \new{קישור לאזור תמונה IIIF.}\\
\hline

\newprop{hm:canonical\_start} & \new{תחילת טווח קנוני.}\\
\hline

\newprop{hm:canonical\_end} & \new{סוף טווח קנוני.}\\
\hline

\newprop{hm:is\_commentary\_on\_canonical} & \new{פירוש על קטעים קנוניים.}\\
\hline
\end{longtable}

%--------------------------------------------------------------------
\subsection{\new{תכונות   — פרדיגמה כפולה}}
%--------------------------------------------------------------------

\begin{longtable}{|r|p{8cm}|}
\hline
\textbf{תכונה} & \textbf{תיאור} \\
\hline
\endhead

\newprop{hm:has\_cataloging\_perspective} & \new{מקשר לתצוגה ביבליוגרפית.}\\
\hline
\newprop{hm:has\_philological\_perspective} & \new{מקשר לתצוגה פילולוגית.}\\
\hline
\newprop{hm:links\_work\_to\_tradition} & \new{מחבר Work ל-TextTradition.}\\
\hline
\newprop{hm:links\_tradition\_to\_work} & \new{מחבר TextTradition ל-Work.}\\
\hline
\newprop{hm:view\_type} & \new{סיווג סוג התצוגה.}\\
\hline
\newprop{hm:cataloging\_work} & \new{Work בתצוגת קטלוג.}\\
\hline
\newprop{hm:philological\_tradition} & \new{מסורת בתצוגה פילולוגית.}\\
\hline
\newprop{hm:is\_paradigm\_compatible} & \new{תצוגות תואמות.}\\
\hline

\newprop{hm:cataloging\_expression} & \new{ביטוי לצרכי קטלוג.}\\
\hline

\newprop{hm:philological\_witness} & \new{עד פילולוגי כאירוע תרבותי.}\\
\hline

\newprop{hm:paradigm\_bridge} & \new{מקשר בין פרדיגמות.}\\
\hline
\end{longtable}

%--------------------------------------------------------------------
\subsection{\new{תכונות   — מבנה CU מקונן}}
%--------------------------------------------------------------------

\begin{longtable}{|r|p{8cm}|}
\hline
\textbf{תכונה} & \textbf{תיאור} \\
\hline
\endhead

\newprop{hm:has\_sub\_unit} & \new{קינון יחידות רקורסיבי (טרנזיטיבי).}\\
\hline
\newprop{hm:is\_sub\_unit\_of} & \new{יחס ילד-הורה.}\\
\hline
\newprop{hm:has\_hierarchy} & \new{מקשר למבנה היררכיה.}\\
\hline
\newprop{hm:hierarchy\_root} & \new{יחידה ברמה עליונה.}\\
\hline
\newprop{hm:hierarchy\_type} & \new{סיווג מורכבות.}\\
\hline
\newprop{hm:precedes\_unit} & \new{סדר רצף פיזי.}\\
\hline

\newprop{hm:follows\_unit} & \new{יחידה עוקבת (הפוך של precedes\_unit).}\\
\hline

\newprop{hm:has\_parent\_unit} & \new{יחידת אב ישירה (לא טרנזיטיבי).}\\
\hline

\newprop{hm:adjacent\_unit} & \new{יחידות סמוכות באותה רמה (סימטרי).}\\
\hline
\end{longtable}

%--------------------------------------------------------------------
\subsection{\new{תכונות   — מסגרת אפיסטמולוגית}}
%--------------------------------------------------------------------

\begin{longtable}{|r|p{8cm}|}
\hline
\textbf{תכונה} & \textbf{תיאור} \\
\hline
\endhead

\newprop{hm:has\_epistemological\_status} & \new{מסווג אופי אפיסטמולוגי.}\\
\hline
\newprop{hm:has\_evidence\_chain} & \new{מקשר טענה לראיות.}\\
\hline
\newprop{hm:evidence\_step} & \new{צעד בודד בשרשרת.}\\
\hline
\newprop{hm:evidence\_source} & \new{מסמך מקור לצעד.}\\
\hline
\newprop{hm:evidence\_agent} & \new{מי סיפק ראיות.}\\
\hline
\newprop{hm:derives\_from} & \new{שרשרת גזירה.}\\
\hline
\newprop{hm:data\_category} & \new{סיווג סוג נתונים.}\\
\hline
\newprop{hm:interpretation\_method} & \new{שיטה לפרשנות.}\\
\hline
\newprop{hm:competing\_interpretation} & \new{ניתוחים חלופיים.}\\
\hline

\newprop{hm:supports\_assertion} & \new{שרשרת ראיות תומכת בטענה.}\\
\hline

\newprop{hm:contradicts\_assertion} & \new{ראיות סותרות טענה.}\\
\hline
\end{longtable}

\newpage

%====================================================================
\section{תכונות נתונים (Datatype Properties)}
%====================================================================

%--------------------------------------------------------------------
\subsection{תכונות CIDOC-CRM}
%--------------------------------------------------------------------

\begin{longtable}{|r|p{5cm}|p{3cm}|}
\hline
\textbf{תכונה} & \textbf{תיאור} & \textbf{סוג} \\
\hline
\endhead

cidoc-crm:P190\_has\_symbolic\_content & תוכן טקסטואלי & xsd:string\\
\hline
cidoc-crm:P82a\_begin\_of\_the\_begin & תחילת חיים )לאדם( & xsd:integer\\
\hline
cidoc-crm:P82b\_end\_of\_the\_end & סוף חיים )לאדם( & xsd:integer\\
\hline
\end{longtable}

%--------------------------------------------------------------------
\subsection{תכונות פיזיות}
%--------------------------------------------------------------------

\begin{longtable}{|r|p{5cm}|p{3cm}|}
\hline
\textbf{תכונה} & \textbf{תיאור} & \textbf{סוג} \\
\hline
\endhead

hm:has\_height\_mm & גובה במ"מ & xsd:decimal\\
\hline
hm:has\_width\_mm & רוחב במ"מ & xsd:decimal\\
\hline
hm:has\_number\_of\_folios & מספר דפים & xsd:integer / xsd:int\\
\hline
hm:has\_lines\_per\_page & שורות בדף & xsd:integer\\
\hline
hm:has\_columns & מספר טורים & xsd:integer\\
\hline
hm:has\_quire\_structure & מבנה קוויונים & xsd:string\\
\hline
hm:has\_catchwords & מילות קשר & xsd:boolean\\
\hline
\end{longtable}

%--------------------------------------------------------------------
\subsection{תכונות טקסטואליות}
%--------------------------------------------------------------------

\begin{longtable}{|r|p{5cm}|p{3cm}|}
\hline
\textbf{תכונה} & \textbf{תיאור} & \textbf{סוג} \\
\hline
\endhead

hm:has\_title & כותרת & xsd:string\\
\hline
hm:has\_incipit & תחילת הטקסט & xsd:string\\
\hline
hm:opening\_text & טקסט פתיחה & xsd:string\\
\hline
hm:has\_explicit & סיום הטקסט & xsd:string\\
\hline
hm:variation\_note & הערת וריאציה & xsd:string\\
\hline
hm:colophon\_text & טקסט הקולופון & xsd:string\\
\hline
hm:has\_folio\_range & טווח דפים & xsd:string\\
\hline

hm:date\_attribution\_source & מקור קביעת התאריך & xsd:string\\
\hline
hm:place\_attribution\_source & מקור קביעת המקום & xsd:string\\
\hline

hm:earliest\_possible\_date & תאריך מוקדם ביותר & xsd:integer\\
\hline

hm:latest\_possible\_date & תאריך מאוחר ביותר & xsd:integer\\
\hline

\new{hm:date\_original\_string} & \new{מחרוזת תאריך מקורית — הטקסט כפי שהופיע במקור} & xsd:string\\
\hline

hm:external\_identifier\_nli & מספר מזהה בספריה הלאומית & xsd:string\\
\hline

hm:external\_uri\_nli & URI מלא לספריה הלאומית & xsd:anyURI\\
\hline

hm:external\_wikidata\_id & מזהה ויקידטה & xsd:string\\
\hline

hm:external\_wikidata\_uri & URI מלא לויקידטה & xsd:anyURI\\
\hline

hm:material\_arrangement & סידור החומר בכתב היד & xsd:string\\
\hline

hm:mentions\_occasion & אירוע המוזכר בקולופון & xsd:string\\
\hline

hm:has\_digital\_representation\_url & קישור לייצוג דיגיטלי & xsd:anyURI\\
\hline

hm:has\_preferred\_citation & אופן ציטוט מועדף & xsd:string\\
\hline

hm:has\_rights\_statement & הצהרת זכויות & xsd:string\\
\hline

hm:has\_ruling\_description & תיאור שרטוט & xsd:string\\
\hline

hm:ownership\_history & היסטוריית בעלות & xsd:string\\
\hline

hm:sfardata\_id & מזהה Sfardata & xsd:string\\
\hline

hm:viaf\_id & מזהה VIAF & xsd:string\\
\hline
\end{longtable}

%--------------------------------------------------------------------
\subsection{\new{תכונות  + — ודאות וייחוס}}
%--------------------------------------------------------------------

\begin{longtable}{|r|p{5cm}|p{3cm}|}
\hline
\textbf{תכונה} & \textbf{תיאור} & \textbf{סוג} \\
\hline
\endhead

\newprop{hm:certainty\_percentage} & \new{אחוז ביטחון} & \new{xsd:decimal}\\
\hline
\newprop{hm:certainty\_note} & \new{הערת ודאות} & \new{xsd:string}\\
\hline
\end{longtable}

%--------------------------------------------------------------------
\subsection{\new{תכונות   — מיקום וגרנולריות}}
%--------------------------------------------------------------------

\begin{longtable}{|r|p{5cm}|p{3cm}|}
\hline
\textbf{תכונה} & \textbf{תיאור} & \textbf{סוג} \\
\hline
\endhead

\newprop{hm:folio\_number} & \new{מספר דף (למשל 15r)} & \new{xsd:string}\\
\hline
\newprop{hm:line\_number} & \new{מספר שורה} & \new{xsd:integer}\\
\hline
\newprop{hm:column\_number} & \new{מספר טור} & \new{xsd:integer}\\
\hline
\newprop{hm:word\_position} & \new{מיקום מילה} & \new{xsd:integer}\\
\hline
\newprop{hm:character\_position} & \new{מיקום תו} & \new{xsd:integer}\\
\hline
\newprop{hm:coordinates\_iiif} & \new{קואורדינטות IIIF} & \new{xsd:string}\\
\hline

\newprop{hm:location\_string} & \new{מחרוזת מיקום קריאה} & \new{xsd:string}\\
\hline

\newprop{hm:variant\_text} & \new{נוסח הגרסה החלופית} & \new{xsd:string}\\
\hline

\newprop{hm:standard\_text} & \new{נוסח מקובל להשוואה} & \new{xsd:string}\\
\hline

\newprop{hm:variant\_significance} & \new{הערכת משמעות הגרסה} & \new{xsd:string}\\
\hline

\newprop{hm:tradition\_name} & \new{שם מסורת טקסט} & \new{xsd:string}\\
\hline

\newprop{hm:tradition\_description} & \new{תיאור מסורת} & \new{xsd:string}\\
\hline

\newprop{hm:is\_foreign\_addition} & \new{האם תוספת זרה (מסנן מהיר)} & \new{xsd:boolean}\\
\hline

\newprop{hm:addition\_period} & \new{תקופת ההוספה} & \new{xsd:string}\\
\hline

\newprop{hm:intervention\_description} & \new{תיאור התערבות סופר} & \new{xsd:string}\\
\hline

\newprop{hm:anthology\_order} & \new{סדר באנתולוגיה} & \new{xsd:integer}\\
\hline

\newprop{hm:overlap\_folio\_range} & \new{טווח דפים חופפים} & \new{xsd:string}\\
\hline

\newprop{hm:volume\_number} & \new{מספר כרך} & \new{xsd:integer}\\
\hline
\end{longtable}

%--------------------------------------------------------------------
\subsection{\new{תכונות   — הפניות קנוניות}}
%--------------------------------------------------------------------

\begin{longtable}{|r|p{5cm}|p{3cm}|}
\hline
\textbf{תכונה} & \textbf{תיאור} & \textbf{סוג} \\
\hline
\endhead

\newprop{hm:book\_name} & \new{שם ספר} & \new{xsd:string}\\
\hline
\newprop{hm:chapter\_number} & \new{מספר פרק} & \new{xsd:integer}\\
\hline
\newprop{hm:verse\_number} & \new{מספר פסוק} & \new{xsd:integer}\\
\hline
\newprop{hm:tractate\_name} & \new{שם מסכת} & \new{xsd:string}\\
\hline
\newprop{hm:mishnah\_number} & \new{מספר משנה} & \new{xsd:integer}\\
\hline
\newprop{hm:talmud\_folio} & \new{דף תלמוד} & \new{xsd:string}\\
\hline
\newprop{hm:halacha\_section} & \new{קטע הלכתי} & \new{xsd:string}\\
\hline
\end{longtable}

%--------------------------------------------------------------------
\subsection{\new{תכונות   — פרדיגמה ומבנה}}
%--------------------------------------------------------------------

\begin{longtable}{|r|p{5cm}|p{3cm}|}
\hline
\textbf{תכונה} & \textbf{תיאור} & \textbf{סוג} \\
\hline
\endhead

\newprop{hm:paradigm\_note} & \new{הערת פרדיגמה} & \new{xsd:string}\\
\hline
\newprop{hm:is\_primary\_paradigm} & \new{פרדיגמה דומיננטית} & \new{xsd:boolean}\\
\hline
\newprop{hm:paradigm\_justification} & \new{הצדקת קישור} & \new{xsd:string}\\
\hline
\newprop{hm:nesting\_level} & \new{רמת קינון} & \new{xsd:integer}\\
\hline
\newprop{hm:is\_atomic\_unit} & \new{יחידה אטומית} & \new{xsd:boolean}\\
\hline
\newprop{hm:max\_nesting\_depth} & \new{עומק קינון מרבי} & \new{xsd:integer}\\
\hline
\end{longtable}

%--------------------------------------------------------------------
\subsection{\new{תכונות   — מסגרת אפיסטמולוגית}}
%--------------------------------------------------------------------

\begin{longtable}{|r|p{5cm}|p{3cm}|}
\hline
\textbf{תכונה} & \textbf{תיאור} & \textbf{סוג} \\
\hline
\endhead

\newprop{hm:is\_factual} & \new{האם עובדה} & \new{xsd:boolean}\\
\hline
\newprop{hm:interpretation\_date} & \new{תאריך פרשנות} & \new{xsd:date}\\
\hline
\newprop{hm:evidence\_strength} & \new{עוצמת ראיה (0-1)} & \new{xsd:decimal}\\
\hline
\newprop{hm:reasoning\_text} & \new{טקסט הנמקה} & \new{xsd:string}\\
\hline
\newprop{hm:methodology\_reference} & \new{הפניה מתודולוגית} & \new{xsd:string}\\
\hline
\newprop{hm:consensus\_level} & \new{רמת הסכמה} & \new{xsd:string}\\
\hline

\newprop{hm:supersedes\_interpretation} & \new{מחליף פרשנות קודמת} & \new{xsd:boolean}\\
\hline

\newprop{hm:unit\_sequence} & \new{מספר סידורי בין יחידות אחיות} & \new{xsd:integer}\\
\hline

\newprop{hm:hierarchy\_description} & \new{תיאור מורכבות היררכיה} & \new{xsd:string}\\
\hline
\end{longtable}

\newpage

%====================================================================
\section{אוצר מילים מבוקר (Instances)}
%====================================================================

%--------------------------------------------------------------------
\subsection{סוגי כתב}
%--------------------------------------------------------------------

\textbf{סוגים אזוריים (TypeScriptType):}
\begin{itemize}
\item hm:AshkenaziScript — כתב אשכנזי
\item hm:SepharadicScript — כתב ספרדי
\item hm:ItalianScript — כתב איטלקי
\item hm:ByzantineScript — כתב ביזנטי
\item hm:YemeniteScript — כתב תימני
\item hm:OrientalScript — כתב מזרחי
\item hm:PersianScript — כתב פרסי
\end{itemize}

\textbf{מצבי כתב (ModeScriptType):}
\begin{itemize}
\item hm:SquareScript — כתב מרובע
\item hm:SemicursiveScript — כתב בינוני
\item hm:CursiveScript — כתב רהוט
\end{itemize}

%--------------------------------------------------------------------
\subsection{חומרים ומצב}
%--------------------------------------------------------------------

\textbf{חומרים:}
\begin{itemize}
\item hm:Parchment — קלף
\item hm:Paper — נייר
\item hm:Papyrus — פפירוס
\end{itemize}

\textbf{מצב פיזי:}
\begin{itemize}
\item hm:Excellent — מצוין
\item hm:Good — טוב
\item hm:Fair — בינוני
\item hm:Poor — גרוע
\item hm:Fragmentary — פרגמנטרי
\end{itemize}

%--------------------------------------------------------------------
\subsection{שפות}
%--------------------------------------------------------------------

\begin{itemize}
\item hm:Hebrew — עברית
\item hm:Aramaic — ארמית
\item hm:Arabic — ערבית
\item hm:JudeoArabic — יהודית-ערבית
\item hm:Ladino — לאדינו
\item hm:Yiddish — יידיש
\item hm:JudeoPersian — יהודית-פרסית
\end{itemize}

%--------------------------------------------------------------------
\subsection{סוגי עיטור}
%--------------------------------------------------------------------

\begin{itemize}
\item hm:FullPageIllumination — עיטור דף שלם
\item hm:InitialLetters — אותיות פותחות מעוטרות
\item hm:MarginalDecoration — עיטור שוליים
\item hm:CarpetPage — דף שטיח
\item hm:Micrography — מיקרוגרפיה
\end{itemize}

%--------------------------------------------------------------------
\subsection{סוגי ניקוד}
%--------------------------------------------------------------------

\begin{itemize}
\item hm:TiberianVocalization — ניקוד טברייני
\item hm:BabylonianVocalization — ניקוד בבלי
\item hm:PalestinianVocalization — ניקוד ארץ-ישראלי
\item hm:NoVocalization — ללא ניקוד
\end{itemize}

%--------------------------------------------------------------------
\subsection{סוגי כריכה}
%--------------------------------------------------------------------

\begin{itemize}
\item hm:OriginalBinding — כריכה מקורית
\item hm:Rebinding — כריכה מחדש
\item hm:ConservationBinding — כריכת שימור
\item hm:Unbound — ללא כריכה
\end{itemize}

%--------------------------------------------------------------------
\subsection{\new{סוגי פורמט תאריך}}
%--------------------------------------------------------------------

\begin{itemize}
\item \new{hm:GregorianYear — שנה גרגוריאנית (למשל: 1407, 1523)}
\item \new{hm:HebrewYear — שנה עברית (למשל: תשפ״ד, ה'קס״ז)}
\item \new{hm:FullDate — תאריך מלא (למשל: 15/03/1407, כ״ג אדר תשפ״ד)}
\item \new{hm:UnstructuredDate — טקסט חופשי (למשל: "המאה ה-15")}
\end{itemize}

%--------------------------------------------------------------------
\subsection{תפקידי השתתפות}
%--------------------------------------------------------------------

\begin{itemize}
\item hm:Scribe\_role — סופר, מעתיק כתב יד
\item hm:Author\_role — מחבר
\item hm:Illuminator\_role — מאייר
\item hm:Commentator\_role — פרשן
\item hm:Translator\_role — מתרגם
\item hm:Owner\_role — בעלים קודם
\end{itemize}

%--------------------------------------------------------------------
\subsection{נושאים}
%--------------------------------------------------------------------

\begin{itemize}
\item hm:Bible\_Subject — מקרא
\item hm:Talmud\_Subject — תלמוד
\item hm:Midrash\_Subject — מדרש
\item hm:Halakha\_Subject — הלכה
\item hm:Kabbalah\_Subject — קבלה
\item hm:Philosophy\_Subject — פילוסופיה
\item hm:Liturgy\_Subject — ליטורגיה
\item hm:Poetry\_Subject — שירה
\item hm:Medicine\_Subject — רפואה
\item hm:Science\_Subject — מדעים
\end{itemize}

%--------------------------------------------------------------------
\subsection{\new{סוגי התערבות סופר}}
%--------------------------------------------------------------------

\begin{itemize}
\item \new{hm:Correction\_type — תיקון}
\item \new{hm:Erasure\_type — מחיקה}
\item \new{hm:Interlinear\_addition\_type — הוספה בין השורות}
\item \new{hm:Marginal\_gloss\_type — ביאור בשוליים}
\item \new{hm:Later\_hand\_type — יד מאוחרת}
\end{itemize}

%--------------------------------------------------------------------
\subsection{\new{סוגי סטטוס יחידה}}
%--------------------------------------------------------------------

\begin{itemize}
\item \new{hm:CoreUnit\_status — יחידת ליבה}
\item \new{hm:LaterAddition\_status — תוספת מאוחרת}
\item \new{hm:BinderAddition\_status — תוספת כורך}
\item \new{hm:UnrelatedFragment\_status — שבר לא קשור}
\item \new{hm:ProtectiveLeaf\_status — דף מגן}
\end{itemize}

%--------------------------------------------------------------------
\subsection{\new{סוגי משמעות גרסה}}
%--------------------------------------------------------------------

\begin{itemize}
\item \new{hm:Orthographic\_variant — גרסה כתיבית (איות)}
\item \new{hm:Lexical\_variant — גרסה לקסיקלית (בחירת מילים)}
\item \new{hm:Syntactic\_variant — גרסה תחבירית}
\item \new{hm:Semantic\_variant — גרסה משמעותית}
\item \new{hm:Addition\_variant — הוספה}
\item \new{hm:Omission\_variant — השמטה}
\end{itemize}

%--------------------------------------------------------------------
\subsection{\new{סוגי היררכיה קודיקולוגית}}
%--------------------------------------------------------------------

\begin{itemize}
\item \new{hm:SimpleHierarchy — היררכיה פשוטה}
\item \new{hm:NestedHierarchy — היררכיה מקוננת}
\item \new{hm:ComplexHierarchy — היררכיה מורכבת (MTM, Sammelband)}
\item \new{hm:FragmentaryHierarchy — היררכיה קטעית}
\end{itemize}

%--------------------------------------------------------------------
\subsection{\new{סוגי פרדיגמה}}
%--------------------------------------------------------------------

\begin{itemize}
\item \new{hm:BibliographicParadigm — פרדיגמה ביבליוגרפית}
\item \new{hm:PhilologicalParadigm — פרדיגמה פילולוגית}
\item \new{hm:HybridParadigm — פרדיגמה היברידית}
\end{itemize}

%--------------------------------------------------------------------
\subsection{סוגי טקסט}
%--------------------------------------------------------------------

\begin{itemize}
\item hm:MainText — טקסט עיקרי
\item hm:CommentaryText — פירוש
\item hm:AdditionalText — טקסט נוסף
\item hm:MarginalText — טקסט בשוליים
\item hm:SupercommentaryText — פירוש על פירוש
\item hm:GlossText — ביאור
\item hm:ResponsumText — תשובה (שו"ת)
\end{itemize}

%--------------------------------------------------------------------
\subsection{\new{רמות ודאות}}
%--------------------------------------------------------------------

\begin{itemize}
\item \new{hm:Certain — ודאי}
\item \new{hm:Probable — סביר}
\item \new{hm:Possible — אפשרי}
\item \new{hm:Uncertain — לא ודאי}
\end{itemize}

%--------------------------------------------------------------------
\subsection{\new{מקורות ייחוס}}
%--------------------------------------------------------------------

\begin{itemize}
\item \new{hm:ExpertAttribution — ייחוס מומחה}
\item \new{hm:AIAttribution — ייחוס AI}
\item \new{hm:CatalogAttribution — ייחוס קטלוג}
\item \new{hm:ColophonAttribution — ייחוס מקולופון}
\item \new{hm:PaleographicAttribution — ייחוס פליאוגרפי}
\end{itemize}

%--------------------------------------------------------------------
\subsection{\new{היררכיות קנוניות}}
%--------------------------------------------------------------------

\begin{itemize}
\item \new{hm:Bible\_hierarchy — היררכיית מקרא}
\item \new{hm:Mishnah\_hierarchy — היררכיית משנה}
\item \new{hm:Talmud\_Bavli\_hierarchy — היררכיית תלמוד בבלי}
\item \new{hm:Talmud\_Yerushalmi\_hierarchy — היררכיית תלמוד ירושלמי}
\item \new{hm:Mishneh\_Torah\_hierarchy — היררכיית משנה תורה}
\item \new{hm:Shulchan\_Aruch\_hierarchy — היררכיית שולחן ערוך}
\item \new{hm:Zohar\_hierarchy — היררכיית ספרות הזוהר}
\item \new{hm:Midrash\_hierarchy — היררכיית מדרש}
\end{itemize}

%--------------------------------------------------------------------
\subsection{\new{סטטוסים אפיסטמולוגיים}}
%--------------------------------------------------------------------

\begin{itemize}
\item \new{hm:DirectObservation — תצפית ישירה}
\item \new{hm:ScholarlyInterpretation — פרשנות חוקר}
\item \new{hm:ComputationalDerivation — גזירה חישובית}
\item \new{hm:CatalogInherited — ירושה מקטלוג}
\item \new{hm:TraditionalAttribution — ייחוס מסורתי}
\end{itemize}

%--------------------------------------------------------------------
\subsection{\new{קטגוריות נתונים}}
%--------------------------------------------------------------------

\begin{itemize}
\item \new{hm:PhysicalMeasurement — מדידה פיזית}
\item \new{hm:ColophonContent — תוכן קולופון}
\item \new{hm:MaterialObservation — תצפית חומר}
\item \new{hm:DateAttribution — ייחוס תאריך}
\item \new{hm:AuthorshipAttribution — ייחוס מחברות}
\item \new{hm:ScribeIdentification — זיהוי סופר}
\item \new{hm:PlaceAttribution — ייחוס מקום}
\item \new{hm:WorkIdentification — זיהוי יצירה}
\item \new{hm:TextualRelationship — יחס טקסטואלי}
\end{itemize}

%--------------------------------------------------------------------
\subsection{\new{שיטות פרשנות}}
%--------------------------------------------------------------------

\begin{itemize}
\item \new{hm:PaleographicAnalysis — ניתוח פליאוגרפי}
\item \new{hm:CodicologicalAnalysis — ניתוח קודיקולוגי}
\item \new{hm:LinguisticAnalysis — ניתוח לשוני}
\item \new{hm:HistoricalContextAnalysis — ניתוח הקשר היסטורי}
\item \new{hm:ComparativeTextualAnalysis — ניתוח טקסטואלי השוואתי}
\item \new{hm:AIBasedAnalysis — ניתוח מבוסס AI}
\item \new{hm:RadiocarbonDating — תיארוך פחמן-14}
\item \new{hm:WatermarkAnalysis — ניתוח סימני מים}
\end{itemize}

\newpage

%====================================================================
\section{עקרונות מידול}
%====================================================================

\subsection{אירועיות}

כל שינוי או פעולה מתועדים כאירוע עם:
\begin{itemize}
\item מבצע \textenglish{(P14\_carried\_out\_by)}
\item זמן \textenglish{(P4\_has\_time-span)}
\item מקום \textenglish{(P7\_took\_place\_at)}
\item חפצים מעורבים \textenglish{(P16\_used\_specific\_object)}
\end{itemize}

\subsection{הפרדת רבדים}

\begin{itemize}
\item \textbf{Work}: התוכן המופשט )רעיון החיבור(
\item \textbf{Expression}: הביטוי הטקסטואלי הספציפי
\item \textbf{Manifestation}: המופע הפיזי )כתב היד(
\end{itemize}

\subsection{גמישות מבנית}

\begin{itemize}
\item כתב יד פשוט: CU יחיד
\item כתב יד מורכב: BU עם מספר CUs
\item יחידות פלאוגרפיות: PU לפי הצורך
\end{itemize}

\subsection{קישוריות}

\begin{itemize}
\item מזהים חיצוניים \textenglish{(VIAF, Wikidata, NLI)}
\item קישור לקטלוגים ומקורות
\item יחסים עם פריטים קשורים
\end{itemize}

\subsection{הרחיבות}

המודל מאפשר הוספת:
\begin{itemize}
\item מחלקות Type חדשות
\item מאפייני נתונים נוספים
\item יחסי עצמים מותאמים
\item אירועים מיוחדים
\end{itemize}

\subsection{השימוש בפועל}

האונטולוגיה מאפשרת שאילתות מורכבות כגון:
\begin{itemize}
\item כל כתבי היד של מחבר מסוים
\item כתבי יד שנוצרו במקום/זמן ספציפיים
\item שרשרת בעלות של כתב יד
\item כתבי יד הכוללים יצירה מסוימת
\item התפוצה הגיאוגרפית של חיבור
\end{itemize}

המודל תומך גם בניתוח השוואתי ובמחקר כמותי על פני אוספים שונים.

\newpage

%====================================================================
\section{דוגמאות RDF}
%====================================================================

%--------------------------------------------------------------------
\subsection{כתב יד פשוט}
%--------------------------------------------------------------------

\begin{LTR}
\begin{verbatim}
# Codicological Unit (single)
:MS_Barberini_82 rdf:type hm:Codicological_Unit ;
    cidoc-crm:P1_is_identified_by :NLI_990000623390205171 ;
    hm:has_number_of_folios 151 ;
    hm:has_material :Paper, :Parchment ;
    hm:has_script_type :SepharadicScript ;
    hm:has_script_mode :SemicursiveScript ;
    hm:has_colophon :MS_Barberini_82_Colophon ;
    lrmoo:R4_embodies :Ginat_Egoz_Expression_MS_Barberini_82 .

# Expression
:Ginat_Egoz_Expression_MS_Barberini_82 rdf:type lrmoo:F2_Expression ;
    lrmoo:R3_is_realised_in :Ginat_Egoz_Work ;
    cidoc-crm:P2_has_type :MainText .

# Work
:Ginat_Egoz_Work rdf:type lrmoo:F1_Work ;
    cidoc-crm:P102_has_title :Ginat_Egoz_Title ;
    cidoc-crm:P2_has_type :Kabbalah_Subject ;
    cidoc-crm:P72_has_language :Hebrew .
\end{verbatim}
\end{LTR}

%--------------------------------------------------------------------
\subsection{כתב יד מורכב \textenglish{(BU + CUs)}}
%--------------------------------------------------------------------

\begin{LTR}
\begin{verbatim}
# Bibliographic Unit (complete volume)
:MS_Vatican_Ebr_44 rdf:type hm:Bibliographic_Unit ;
    cidoc-crm:P1_is_identified_by :MS_Vatican_Ebr_44_Identifier ;
    hm:has_number_of_folios 284 ;
    lrmoo:R5_has_component :MS_Vatican_Ebr_44_CU_01, 
                           :MS_Vatican_Ebr_44_CU_02 .

# Codicological Unit 1
:MS_Vatican_Ebr_44_CU_01 rdf:type hm:Codicological_Unit ;
    lrmoo:R5i_is_component_of :MS_Vatican_Ebr_44 ;
    lrmoo:R4_embodies :Midrash_Tanchuma_Expression_MS_Vatican_44 .

# Codicological Unit 2
:MS_Vatican_Ebr_44_CU_02 rdf:type hm:Codicological_Unit ;
    lrmoo:R5i_is_component_of :MS_Vatican_Ebr_44 ;
    lrmoo:R4_embodies :Alphabet_Ben_Sira_Expression_MS_Vatican_44 .
\end{verbatim}
\end{LTR}

%--------------------------------------------------------------------
\subsection{אירוע הפקה וקולופון}
%--------------------------------------------------------------------

\begin{LTR}
\begin{verbatim}
# Production Event
:MS_Barberini_82_Production_Event rdf:type lrmoo:F32_Item_Production_Event ;
    lrmoo:R27_materialized :MS_Barberini_82 ;
    cidoc-crm:P14_carried_out_by :Moshe_ben_Yitzhak_Ibn_Tibbon ;
    cidoc-crm:P4_has_time-span :Year_1407_Shvat ;
    cidoc-crm:P7_took_place_at :Candia_Herakleion .

:Year_1407_Shvat rdf:type cidoc-crm:E52_Time-Span .
:Moshe_ben_Yitzhak_Ibn_Tibbon rdf:type cidoc-crm:E21_Person .
:Candia_Herakleion rdf:type cidoc-crm:E53_Place .

# Colophon
:MS_Barberini_82_Colophon rdf:type hm:Colophon ;
    hm:mentions_scribe :Moshe_ben_Yitzhak_Ibn_Tibbon ;
    hm:mentions_date :Year_1407_Shvat ;
    hm:mentions_place :Candia_Herakleion ;
    hm:colophon_text "חק אבן תבון וסיימתי בחדש שבט שנת הקס\"ז לבריאת עולם למינן שאנו מונין פה בעיר קנדיאה"@he .
\end{verbatim}
\end{LTR}

%--------------------------------------------------------------------
\subsection{אירוע העברת בעלות}
%--------------------------------------------------------------------

\begin{LTR}
\begin{verbatim}
:MS_Vatican_Ebr_44_Ownership_History_Event rdf:type cidoc-crm:E10_Transfer_of_Custody ;
    cidoc-crm:P16_used_specific_object :MS_Vatican_Ebr_44 ;
    cidoc-crm:P24_transferred_title_from :Elijah_ben_Elkanah_Capsali ;
    cidoc-crm:P23_transferred_title_to :Anton_Fugger ;
    cidoc-crm:P4_has_time-span :November_8_1541 .

:November_8_1541 rdf:type cidoc-crm:E52_Time-Span .
\end{verbatim}
\end{LTR}

\newpage

%====================================================================
\section{סיכום השינויים מהעיצוב המקורי}
%====================================================================

האונטולוגיה המעודכנת מרחיבה את העיצוב המקורי של פרופ' גילה פרבור במספר כיוונים מרכזיים:

\subsection{תמיכה בפרדיגמה כפולה}

\new{נוספה תמיכה בשתי גישות שונות לתיאור כתבי יד:}
\begin{itemize}
\item \new{\textbf{גישה ביבליוגרפית-מעשית:} מבוססת על מודל Work-Expression, מתאימה לקטלוג וויקידטה.}
\item \new{\textbf{גישה פילולוגית חדשה:} מתייחסת לכל כתב יד כאירוע תרבותי ייחודי, ללא הנחת "יצירה" אחידה.}
\end{itemize}

\subsection{מבנה קודיקולוגי מקונן}

\new{הותר קינון של יחידות קודיקולוגיות בעומק שרירותי:}
\begin{itemize}
\item \new{CU יכולה להכיל תת-CU שמכילה תת-תת-CU וכו'.}
\item \new{מאפשר ייצוג כתבי יד מורכבים כמו Sammelbands ו-MTM.}
\end{itemize}

\subsection{מסגרת אפיסטמולוגית}

\new{נוספה הבחנה שיטתית בין סוגי נתונים:}
\begin{itemize}
\item \new{\textbf{תצפיות ישירות:} עובדות ניתנות לאימות (מידות, ספירת דפים).}
\item \new{\textbf{פרשנויות חוקרים:} הכרעות הדורשות שיפוט מומחה.}
\item \new{\textbf{גזירות חישוביות:} תוצאות אלגוריתמים ו-AI.}
\end{itemize}

\subsection{תמיכה בגרנולריות מלאה}

\new{נוספה תמיכה בתיאור ברמות שונות:}
\begin{itemize}
\item \new{דף, שורה, מילה, תו}
\item \new{אינטגרציה עם IIIF לאזורי תמונה}
\end{itemize}

\subsection{היררכיות קנוניות}

\new{נוספה תמיכה בשמונה היררכיות טקסט קנוניות יהודיות:}
\begin{itemize}
\item \new{מקרא (ספר/פרק/פסוק)}
\item \new{משנה (מסכת/פרק/משנה)}
\item \new{תלמוד בבלי וירושלמי (מסכת/דף/עמוד)}
\item \new{משנה תורה ושולחן ערוך (קטע הלכתי)}
\item \new{זוהר ומדרש (מבנה ספרותי ייחודי)}
\end{itemize}

\newpage

%====================================================================
\section{נספח: טבלאות מהתיעוד 26.09.2025}
%====================================================================

נספח זה משחזר את טבלאות המחלקות והמאפיינים כפי שהופיעו במסמך \textenglish{``תיעוד אונטולוגיה לכתבי יד עבריים - גרסה מעודכנת (3)''} מתאריך 26.09.2025.

%--------------------------------------------------------------------
\subsection{מחלקות (Classes) — כפי שמופיע בתיעוד 26-9-25}
%--------------------------------------------------------------------

{\small
\begin{longtable}{|p{4.5cm}|p{8.5cm}|p{3cm}|}
\hline
\textbf{מחלקה} & \textbf{תיאור} & \textbf{סופרקלאס} \\
\hline
\endhead

\textenglish{lrmoo:F1\_Work} & יצירה מופשטת )רעיון התוכן(. & — \\
\hline
\textenglish{lrmoo:F2\_Expression} & ביטוי טקסטואלי ספציפי. & — \\
\hline
\textenglish{lrmoo:F3\_Manifestation} & מהדורה כללית. & \textenglish{cidoc-crm:E22\_Human-Made\_Object} \\
\hline
\textenglish{lrmoo:F4\_Manifestation\_Singleton} & עותק פיזי יחידאי )כתב יד(. & \textenglish{lrmoo:F3\_Manifestation} \\
\hline
\textenglish{lrmoo:F5\_Item} & פריט פיזי \textenglish{(LRMoo)} תואם. & \textenglish{cidoc-crm:E18\_Physical\_Thing} \\
\hline
\end{longtable}
}

{\small
\begin{longtable}{|p{4.5cm}|p{8.5cm}|p{3cm}|}
\hline
\textbf{מחלקה} & \textbf{תיאור} & \textbf{יוצר/קשר} \\
\hline
\endhead

\textenglish{lrmoo:F27\_Work\_Creation} & אירוע יצירה של \textenglish{Work}. & \textenglish{R16\_created → F1\_Work} \\
\hline
\textenglish{lrmoo:F28\_Expression\_Creation} & אירוע יצירה של \textenglish{Expression}. & \textenglish{R17\_created → F2\_Expression} \\
\hline
\textenglish{lrmoo:F30\_Manifestation\_Creation} & אירוע יצירה של \textenglish{Manifestation}. & \textenglish{R24\_created → F3\_Manifestation} \\
\hline
\textenglish{lrmoo:F32\_Item\_Production\_Event} & אירוע הפקת עותק )העתקה(. & \textenglish{R27\_materialized → F4\_Manifestation\_Singleton} \\
\hline
\end{longtable}
}

{\small
\begin{longtable}{|p{4.5cm}|p{8.5cm}|p{3cm}|}
\hline
\textbf{מחלקה} & \textbf{תיאור} & \textbf{סופרקלאס} \\
\hline
\endhead

\textenglish{cidoc-crm:E7\_Activity} & פעילות כללית. & \textenglish{cidoc-crm:E1\_CRM\_Entity} \\
\hline
\textenglish{cidoc-crm:E12\_Production} & יצירת חפץ פיזי. & \textenglish{cidoc-crm:E7\_Activity} \\
\hline
\textenglish{cidoc-crm:E8\_Acquisition} & רכישה. & \textenglish{cidoc-crm:E7\_Activity} \\
\hline
\textenglish{cidoc-crm:E10\_Transfer\_of\_Custody} & העברת משמורת. & \textenglish{cidoc-crm:E7\_Activity} \\
\hline
\textenglish{cidoc-crm:E21\_Person} & אדם. & \textenglish{cidoc-crm:E39\_Actor} \\
\hline
\textenglish{cidoc-crm:E74\_Group} & ארגון/קבוצה. & \textenglish{cidoc-crm:E39\_Actor} \\
\hline
\textenglish{cidoc-crm:E52\_Time-Span} & פרק זמן. & \textenglish{cidoc-crm:E1\_CRM\_Entity} \\
\hline
\textenglish{cidoc-crm:E53\_Place} & מקום. & \textenglish{cidoc-crm:E1\_CRM\_Entity} \\
\hline
\textenglish{cidoc-crm:E55\_Type} & טיפוס/סוג. & — \\
\hline
\textenglish{cidoc-crm:E56\_Language} & שפה. & \textenglish{cidoc-crm:E55\_Type} \\
\hline
\textenglish{cidoc-crm:E57\_Material} & חומר. & \textenglish{cidoc-crm:E55\_Type} \\
\hline
\end{longtable}
}

{\small
\begin{longtable}{|p{4.5cm}|p{8.5cm}|p{3cm}|}
\hline
\textbf{מחלקה} & \textbf{תיאור} & \textbf{סופרקלאס} \\
\hline
\endhead

\textenglish{HebrewManuscripts:Bibliographic\_Unit} & יחידה ביבליוגרפית )כרך שלם(. & \textenglish{F4\_Manifestation\_Singleton} \\
\hline
\textenglish{HebrewManuscripts:Codicological\_Unit} & יחידה קודיקולוגית )חיבור בזמן/במקום(. & \textenglish{F4\_Manifestation\_Singleton} \\
\hline
\textenglish{HebrewManuscripts:Paleographical\_Unit} & יחידה פלאוגרפית )יד כותב(. & \textenglish{F4\_Manifestation\_Singleton} \\
\hline
\textenglish{HebrewManuscripts:Colophon} & קולופון )סיום כתב יד(. & \textenglish{cidoc-crm:E73\_Information\_Object} \\
\hline
\textenglish{HebrewManuscripts:Reference\_Event} & אזכור/אירוע הפניה במחקר. & \textenglish{cidoc-crm:E7\_Activity} \\
\hline
\end{longtable}
}

{\small
\begin{longtable}{|p{4.5cm}|p{8.5cm}|p{3cm}|}
\hline
\textbf{מחלקה} & \textbf{תיאור} & \textbf{סופרקלאס} \\
\hline
\endhead

\textenglish{HebrewManuscripts:HebrewScriptType} & סוג כתב עברי. & \textenglish{cidoc-crm:E55\_Type} \\
\hline
\textenglish{HebrewManuscripts:TypeScriptType} & סוג כתב לפי מסורת )ספרדי/אשכנזי(. & \textenglish{HebrewScriptType} \\
\hline
\textenglish{HebrewManuscripts:ModeScriptType} & סגנון כתב )מרובע/קורסיבי(. & \textenglish{HebrewScriptType} \\
\hline
\textenglish{HebrewManuscripts:SubjectType} & נושא תוכן )קבלה, הלכה וכו'(. & \textenglish{cidoc-crm:E55\_Type} \\
\hline
\textenglish{HebrewManuscripts:ParticipationRole} & תפקיד השתתפות. & \textenglish{cidoc-crm:E55\_Type} \\
\hline
\end{longtable}
}

%--------------------------------------------------------------------
\subsection{יחסי עצמים (Object Properties) — כפי שמופיע בתיעוד 26-9-25}
%--------------------------------------------------------------------

{\small
\begin{longtable}{|p{4cm}|p{3.5cm}|p{3.5cm}|p{4.5cm}|}
\hline
\textbf{יחס} & \textbf{תחום} & \textbf{טווח} & \textbf{תיאור} \\
\hline
\endhead

\textenglish{R3\_is\_realised\_in} & \textenglish{F1\_Work} & \textenglish{F2\_Expression} & יצירה מתממשת בביטוי \\
\hline
\textenglish{R4i\_is\_embodied\_in} & \textenglish{F2\_Expression} & \textenglish{F3/F4\_Manifestation} & ביטוי מגולם במהדורה \\
\hline
\textenglish{R7\_exemplifies} & \textenglish{F4\_Manifestation\_Singleton} & \textenglish{F3\_Manifestation} & עותק מגלם מהדורה \\
\hline
\textenglish{R10\_has\_member} & \textenglish{F1\_Work} & \textenglish{F1\_Work} & בת-יצירה כוללת יצירות \\
\hline
\textenglish{R5\_has\_component} & \textenglish{F4\_Manifestation\_Singleton} & \textenglish{F3/F4\_Manifestation} & הרכבת חלקים \textenglish{(BU→CU→PU)} \\
\hline
\textenglish{R69\_has\_physical\_form} & \textenglish{F3/F4\_Manifestation} & \textenglish{E55\_Type} & צורה פיזית \\
\hline
\textenglish{R16\_created} & \textenglish{F27\_Work\_Creation} & \textenglish{F1\_Work} & יצירת \textenglish{Work} \\
\hline
\textenglish{R17\_created} & \textenglish{F28\_Expression\_Creation} & \textenglish{F2\_Expression} & יצירת \textenglish{Expression} \\
\hline
\textenglish{R24\_created} & \textenglish{F30\_Manifestation\_Creation} & \textenglish{F3\_Manifestation} & יצירת \textenglish{Manifestation} \\
\hline
\textenglish{R27\_materialized} & \textenglish{F32\_Item\_Production\_Event} & \textenglish{F4\_Manifestation\_Singleton} & הפקת עותק פיזי \\
\hline
\textenglish{P1\_is\_identified\_by} & \textenglish{F4\_Manifestation\_Singleton} & \textenglish{E42\_Identifier} & מזהה כתב יד \\
\hline
\textenglish{P2\_has\_type} & כללי & \textenglish{E55\_Type} & שיוך לטיפוס \\
\hline
\textenglish{P72\_has\_language} & \textenglish{F2\_Expression} & \textenglish{E56\_Language} & שפת הטקסט \\
\hline
\textenglish{P4\_has\_time-span} & \textenglish{E7\_Activity} & \textenglish{E52\_Time-Span} & זמן אירוע \\
\hline
\textenglish{P7\_took\_place\_at} & \textenglish{E7\_Activity} & \textenglish{E53\_Place} & מקום אירוע \\
\hline
\textenglish{P14\_carried\_out\_by} & \textenglish{E7\_Activity} & \textenglish{E39\_Actor} & מבצע האירוע \\
\hline
\textenglish{P16\_used\_specific\_object} & \textenglish{E7\_Activity} & \textenglish{E18\_Physical\_Thing} & שימוש בחפץ \\
\hline
\textenglish{P23\_transferred\_title\_to} & \textenglish{E8\_Acquisition} & \textenglish{E39\_Actor} & בעלים חדש \\
\hline
\textenglish{P24\_transferred\_title\_from} & \textenglish{E8\_Acquisition} & \textenglish{E39\_Actor} & בעלים קודם \\
\hline
\textenglish{P28\_custody\_surrendered\_by} & \textenglish{E10\_Transfer\_of\_Custody} & \textenglish{E39\_Actor} & מסר משמורת \\
\hline
\textenglish{P29\_custody\_received\_by} & \textenglish{E10\_Transfer\_of\_Custody} & \textenglish{E39\_Actor} & קיבל משמורת \\
\hline
\textenglish{P70i\_is\_documented\_in} & \textenglish{F4\_Manifestation\_Singleton} & \textenglish{F3\_Manifestation} & תועד בקטלוג \\
\hline
\textenglish{P102\_has\_title} & \textenglish{E22\_Human-Made\_Object} & \textenglish{E35\_Title} & כותרת \\
\hline
\textenglish{has\_colophon} & \textenglish{F4\_Manifestation\_Singleton} & \textenglish{Colophon} & קישור לקולופון \\
\hline
\textenglish{has\_script\_type} & \textenglish{F4\_Manifestation\_Singleton} & \textenglish{E55\_Type} & סוג הכתב \\
\hline
\textenglish{has\_Script\_Mode} & \textenglish{F4\_Manifestation\_Singleton} & \textenglish{E55\_Type} & סגנון הכתב \\
\hline
\textenglish{has\_material} & \textenglish{F4\_Manifestation\_Singleton} & \textenglish{E57\_Material} & חומר הכתיבה \\
\hline
\textenglish{has\_date\_of\_creation} & \textenglish{F4\_Manifestation\_Singleton} & \textenglish{E52\_Time-Span} & תאריך יצירה \\
\hline
\textenglish{mentions\_scribe} & \textenglish{Colophon} & \textenglish{E21\_Person} & סופר המוזכר בקולופון \\
\hline
\textenglish{mentions\_date} & \textenglish{Colophon} & \textenglish{E52\_Time-Span} & תאריך בקולופון \\
\hline
\textenglish{mentions\_place} & \textenglish{Colophon} & \textenglish{E53\_Place} & מקום בקולופון \\
\hline
\end{longtable}
}

%--------------------------------------------------------------------
\subsection{מאפייני נתונים (Data Properties) — כפי שמופיע בתיעוד 26-9-25}
%--------------------------------------------------------------------

{\small
\begin{longtable}{|p{4cm}|p{3.5cm}|p{3.5cm}|p{4.5cm}|}
\hline
\textbf{מאפיין} & \textbf{תחום} & \textbf{טווח} & \textbf{תיאור} \\
\hline
\endhead

\textenglish{P190\_has\_symbolic\_content} & \textenglish{E33\_Linguistic\_Object} & \textenglish{xsd:string} & תוכן טקסטואלי \\
\hline
\textenglish{P82a\_begin\_of\_the\_begin} & \textenglish{E21\_Person} & \textenglish{xsd:integer} & תחילת חיים \\
\hline
\textenglish{P82b\_end\_of\_the\_end} & \textenglish{E21\_Person} & \textenglish{xsd:integer} & סוף חיים \\
\hline
\textenglish{colophon\_text} & \textenglish{Colophon} & \textenglish{xsd:string} & טקסט הקולופון \\
\hline
\textenglish{mentions\_occasion} & \textenglish{Colophon} & \textenglish{xsd:string} & אירוע המוזכר \\
\hline
\textenglish{earliest\_possible\_date} & \textenglish{E52\_Time-Span} & \textenglish{xsd:integer} & תאריך מוקדם ביותר \\
\hline
\textenglish{latest\_possible\_date} & \textenglish{E52\_Time-Span} & \textenglish{xsd:integer} & תאריך מאוחר ביותר \\
\hline
\textenglish{date\_attribution\_source} & כללי & \textenglish{xsd:string} & מקור ייחוס התאריך \\
\hline
\textenglish{place\_attribution\_source} & \textenglish{F4\_Manifestation\_Singleton} & \textenglish{xsd:string} & מקור ייחוס המקום \\
\hline
\textenglish{has\_number\_of\_folios} & \textenglish{F4\_Manifestation\_Singleton} & \textenglish{xsd:int} & מספר דפים \\
\hline
\textenglish{has\_folio\_range} & כללי & \textenglish{xsd:string} & טווח דפים \\
\hline
\textenglish{has\_height\_mm} & כללי & \textenglish{xsd:decimal} & גובה במ"מ \\
\hline
\textenglish{has\_width\_mm} & כללי & \textenglish{xsd:decimal} & רוחב במ"מ \\
\hline
\textenglish{material\_arrangement} & \textenglish{F4\_Manifestation\_Singleton} & \textenglish{xsd:string} & סידור חומרים \\
\hline
\textenglish{opening\_text} & \textenglish{F2\_Expression} & \textenglish{xsd:string} & טקסט פתיחה \\
\hline
\textenglish{variation\_note} & \textenglish{F2\_Expression} & \textenglish{xsd:string} & הערת וריאציה \\
\hline
\textenglish{external\_identifier\_nli} & \textenglish{F4\_Manifestation\_Singleton} & \textenglish{xsd:string} & מזהה ספרייה לאומית \\
\hline
\textenglish{external\_uri\_nli} & \textenglish{SubjectType} & \textenglish{xsd:anyURI} & \textenglish{URI} ספרייה לאומית \\
\hline
\textenglish{external\_wikidata\_id} & \textenglish{owl:Thing} & \textenglish{xsd:string} & מזהה \textenglish{Wikidata} \\
\hline
\textenglish{external\_wikidata\_uri} & \textenglish{owl:Thing} & \textenglish{xsd:anyURI} & \textenglish{URI Wikidata} \\
\hline
\textenglish{viaf\_id} & \textenglish{E21\_Person} & \textenglish{xsd:string} & מזהה \textenglish{VIAF} \\
\hline
\textenglish{sfardata\_id} & כללי & \textenglish{xsd:string} & מזהה \textenglish{Sfardata} \\
\hline
\textenglish{has\_digital\_representation\_url} & \textenglish{F4\_Manifestation\_Singleton} & \textenglish{xsd:anyURI} & קישור לדיגיטציה \\
\hline
\textenglish{ownership\_history} & כללי & \textenglish{xsd:string} & היסטוריית בעלות \\
\hline
\textenglish{P44\_has\_condition} & כללי & \textenglish{E55\_Type} & מצב שימור \\
\hline
\end{longtable}
}

\section{\new{הנחיות שימוש ועדכונים \en{(גרסה 1.6)}}
\subsection{מעבר מתכונות שהוצאו משימוש \en{(Deprecated)}}
התכונות הבאות מסומנות \textenglish{owl:deprecated true} ו-\textenglish{rdfs:seeAlso} מצביע למחליף. יש להעדיף את התכונה החדשה בנתונים חדשים:
\begin{itemize}
\item \textenglish{hm:has\_Script\_Mode} $\rightarrow$ \textenglish{hm:has\_script\_mode}
\item \textenglish{hm:external\_wikidata\_id} $\rightarrow$ \textenglish{hm:wikidata\_id}
\item \textenglish{hm:links\_work\_to\_tradition} $\rightarrow$ \textenglish{hm:has\_linked\_work}
\item \textenglish{hm:links\_tradition\_to\_work} $\rightarrow$ \textenglish{hm:has\_linked\_tradition}
\item \textenglish{hm:anthology\_position} \en{(datatype)} $\rightarrow$ \textenglish{hm:has\_anthology\_position} + \textenglish{hm:AnthologyPosition} + \textenglish{hm:anthology\_order}
\end{itemize}

\subsection{דפוס מיקום באנתולוגיה \en{(Anthology Position)}}
הדפוס המומלץ: \textenglish{Expression} $\rightarrow$ \textenglish{AnthologyPosition} $\rightarrow$ \textenglish{xsd:integer} \en{(via has\_anthology\_position then anthology\_order)}. אין להשתמש בתכונה \textenglish{hm:anthology\_position} \en{(Expression $\rightarrow$ integer)} בנתונים חדשים.

\subsection{תכונות הופכיות חדשות}
\begin{itemize}
\item \textenglish{hm:was\_former\_owner\_of} — הופכי ל-\textenglish{hm:former\_owner} \en{(Person $\rightarrow$ Manuscript)}
\item \textenglish{hm:hosts\_part} — הופכי ל-\textenglish{hm:is\_part\_of\_host} \en{(host item $\rightarrow$ part manuscripts)}
\end{itemize}

\subsection{אקסיומת אי-חפיפה \en{(Disjointness)}}
מחלקות תת-סוג של \textenglish{E55\_Type} הבאות מוכרזות \textenglish{owl:AllDisjointClasses}: \textenglish{ConditionType}, \textenglish{VocalizationType}, \textenglish{BindingType}, \textenglish{DateFormatType}, \textenglish{CertaintyLevel}, \textenglish{AttributionSource}, \textenglish{DecorationType}, \textenglish{SubjectType}, \textenglish{ParticipationRole}, \textenglish{CanonicalHierarchyType}, \textenglish{UnitStatusType}. מחלקות המיקום \en{(FolioLocation, LineLocation, WordLocation, CharacterLocation)} כבר מוכרזות כאי-חופפות.

\vspace{2em}
\begin{center}
\rule{0.5\textwidth}{0.4pt}\\[1em]
\textbf{סוף המסמך}
\end{center}



\newpage
% =============================================================================
\section*{סיכום}
% =============================================================================

מסמך זה תיעד את המיפוי המלא בין שדות MARC לישויות ומאפיינים באונטולוגיית כתבי יד עבריים גרסה 1.5. הנקודות העיקריות:

\begin{enumerate}
\item \textbf{מיפוי מקיף}: כל שדות ה-MARC הרלוונטיים ממופים לאונטולוגיה \textenglish{(001, 008, 041, 100, 110, 245, 246, 260/264, 300, 340, 500, 505, 510, 561, 563, 6XX, 700/710, 856)}
\item \textbf{LRM-OO}: שימוש במודל Work-Expression-Manifestation עם תמיכה מלאה ב-F1\_Work, F2\_Expression, F4\_Manifestation\_Singleton
\item \textbf{CIDOC-CRM}: מיפוי מלא לאירועים \textenglish{(E7, E12)}, מקומות \textenglish{(E53)}, זמנים \textenglish{(E52)}, אנשים \textenglish{(E21)} וארגונים \textenglish{(E74)}
\item \textbf{מסגרת אפיסטמולוגית}: הבחנה בין עובדות לפרשנויות, כולל 5 סטטוסים \textenglish{(DirectObservation, ScholarlyInterpretation, ComputationalDerivation, CatalogInherited, TraditionalAttribution)}, 8 שיטות פרשנות, 9 קטגוריות נתונים
\item \textbf{רמות וודאות}: 4 רמות \textenglish{(Certain, Probable, Possible, Uncertain)} עם זיהוי אוטומטי של סימני אי-ודאות מ-MARC
\item \textbf{מקורות ייחוס}: 5 מקורות \textenglish{(CatalogAttribution, ColophonAttribution, PaleographicAttribution, ExpertAttribution, AIAttribution)}
\item \textbf{פרדיגמה כפולה}: תמיכה בגישה ביבליוגרפית \textenglish{(CatalogingView)} ופילולוגית \textenglish{(PhilologicalView)} עם ParadigmBridge, 3 סוגי פרדיגמה \textenglish{(Bibliographic, Philological, Hybrid)}, 4 סוגי היררכיה \textenglish{(Simple, Nested, Complex, Fragmentary)}
\item \textbf{מבנה מקונן}: תמיכה ביחידות קודיקולוגיות מורכבות \textenglish{(CodicologicalHierarchy, CompositeCodicologicalUnit, AtomicCodicologicalUnit)} עם \textenglish{has\_sub\_unit} טרנזיטיבי
\item \textbf{תבניות URI}: 30 תבניות URI ייחודיות לכל סוגי הישויות )כולל v1.4: Person, Group, Place, TimeSpan, Production, WorkCreation, Ownership, Catalog, Subject, Binding, BibUnit, CU, PU, MultiVolumeSet, AnthologyStructure, Hierarchy, CatalogingView, PhilologicalView, TextTradition, Witness, Bridge, Variant, Intervention, CanonRef, Loc, EvidenceChain(
\item \textbf{אוצרות מילים מבוקרים}: 19 שפות, 7 סוגי כתב, 3 מצבי כתב, 8 היררכיות קאנוניות )כולל זוהר ומדרש(, 5 סוגי מצב, 4 סוגי ניקוד, 4 סוגי כריכה, 4 סוגי פורמט תאריך, 5 סוגי קישוט, 12 ז'אנרים, 7 סוגי טקסט )מבני(, 5 סוגי התערבות, 6 סוגי חשיבות גרסה, 5 סוגי סטטוס יחידה
\item \textbf{API מתקדם}: 10 פונקציות API לתכונות v1.4 שאינן נכתבות אוטומטית מ-MARC )כולל philological view, text tradition, paradigm bridge(
\item \textbf{שדות ExtractedData מורחבים}: 12 שדות v1.4 \textenglish{(certainty\_levels, attribution\_sources, data\_from\_colophon, codicological\_units, hierarchy\_type, is\_anthology, is\_multi\_volume, volume\_number, scribal\_interventions, canonical\_references, text\_traditions, textual\_variants)}
\item \textbf{מטא-נתונים מנהליים \textenglish{(v1.5)}}: כיסוי מלא של שדות מנהליים ונתוני גישה מרשומות \textenglish{MARC} של הספרייה הלאומית:
  \begin{itemize}
  \item \textenglish{544} — חומרים קשורים \textenglish{(hm:related\_material\_note)}
  \item \textenglish{541} — מקור רכישה \textenglish{(hm:acquisition\_source)}
  \item \textenglish{597} — הודעת זכויות יוצרים \textenglish{(hm:copyright\_notice)}
  \item \textenglish{939} — הגבלות שימוש \textenglish{(hm:UsageRestriction)}
  \item \textenglish{952} — קביעת זכויות \textenglish{(hm:RightsDetermination)}
  \item \textenglish{966} — שם אוסף \textenglish{(hm:collection\_name)}
  \item \textenglish{AVA} — אחזקות פיזיות \textenglish{(hm:PhysicalHolding)}
  \item \textenglish{AVE} — נקודות גישה דיגיטליות \textenglish{(hm:DigitalAccess)}
  \end{itemize}
\end{enumerate}

\textbf{תאימות:} מסמך זה מסונכרן עם:
\begin{itemize}
\item אונטולוגיה: \textenglish{hebrew-manuscripts.ttl} )v1.5, פברואר 2026(
\item ממיר — קבצי תצורה:
  \begin{itemize}
  \item \textenglish{converter/config/namespaces.py} — 8 מרחבי שמות
  \item \textenglish{converter/config/vocabularies.py} — אוצרות מילים מבוקרים
  \end{itemize}
\item ממיר — לוגיקת המרה:
  \begin{itemize}
  \item \textenglish{converter/transformer/field\_handlers.py} — חילוץ נתונים מ-MARC
  \item \textenglish{converter/transformer/uri\_generator.py} — 30+ תבניות URI
  \item \textenglish{converter/rdf/graph\_builder.py} — בניית גרף RDF עם API מתקדם
  \end{itemize}
\end{itemize}

\textbf{עדכון אחרון:} פברואר 2026

\end{document}

